%% 第26回編集委員会（物理第4巻）：組 電W4
%% p4-06.tex  電-6　電流が方位磁針を動かす：電流がつくる磁場（★1〜4・3段にしない）
%% 中身は計画書 tools/phys4-plan.md の B「電-6」。公式の持ち主は A（A11・A12・A38〜A44・A69〜A71）。
%%   磁場の記号を巻ではじめて出す章（Y-2：E-B 対応）。B を「磁場（磁束密度）」とよび、H は電-7 で「磁場の強さ」として出す。
%%   高校の H と B＝μ0H は③で一言。「電界・磁界」は専門書の囲みで読み方だけ。μ0＝4π×10⁻⁷ と 2019年の SI の改定は③（Y-6）。
%%   電-1〜5 はほかの組が同時に書いているので、計画書 B の節の番号で指した（電-1 ④電場・⑧電場の中の荷電粒子、電-2 ⑥電子ボルト・例題5、
%%   電-4 ①電流・②I＝envS）。電-8（③磁束・⑥自己誘導）、電-12（③直流モーター）、電-14（変位電流）も計画書の番号。
%% 数学との重なり（Y-9）：4-14⑧（ローレンツ力・3 A・0.2 m・0.5 T・0.3 N・フレミングの左手）、4-14④（右ねじ・⊙⊗）、
%%   8-10⑪（B＝μ0I/(2πr)・10 A・5 cm・4×10⁻⁵ T・rot B＝μ0 j）、8-11⑬（アンペールの法則は道を選ぶ）、8-9（div B＝0）は、
%%   同じ数・同じ言い方で指して使った。物理の章では「どの実験・法則から出るか」「実物の値」「仮定がくずれるとき」を主にした。
%% 数値はすべて Python で計算した：checks/物理第4巻/電-6/calc_p4_06.py（出力 calc_p4_06_out.txt）。⑥の囲みのプログラムは
%%   同じ計算（有限のソレノイドの軸の上の磁場）。常用漢字の確かめ：checks/物理第4巻/電-6/kanji_p4_06_07.py
%% 記号（この章）：磁場（磁束密度）B、真空の透磁率 μ0、電流 I、電荷 q（粒子）、電気素量 e、速さ v、粒子の質量 m、
%%   導線から点までの距離 r、円電流の半径 a、巻数 N・1 m あたりの巻数 n、導線の長さ l（数学 4-14 は \vec L）、
%%   トルク T（Y-4。周期はこの章では使わず周波数 f で書く。トルクの記号カードは電-12 に置く：Y-4）、コイルの面積 S、
%%   ホール電圧 V_H、数密度 n_c（巻数の n とぶつけないため、キャリアの数密度は n_c と書いた）、素子の厚さ t_H（時刻 t とぶつけない）、
%%   電場 E（電-1 と同じ）。偏角・伏角は記号を作らず、ことばで書いた（計画書）。磁気モーメントは記号を作らず「NIS」と書いた（C-1）。
%% 出典（書くときに確かめた。11月5日）
%%   ・μ0：CODATA 2022 1.25663706127×10⁻⁶ N/A²（NIST https://physics.nist.gov/cgi-bin/cuu/Value?mu0 、curl で読んだ）。
%%       4π×10⁻⁷ との差は相対 −1.3×10⁻¹⁰（calc_p4_06.py）。
%%   ・アンペアの定義：今の定義（e＝1.602 176 634×10⁻¹⁹ C に固定）BIPM https://www.bipm.org/en/si-base-units/ampere（curl で読んだ）。
%%       前の定義（1948年1月1日から：無限に長い平行な2本、1 m、2×10⁻⁷ N/m）、「前の定義は μ0 を 4π×10⁻⁷ に固定していた。
%%       新しい定義では μ0 は実験で決める」、改定の施行 2019年5月20日、改定のときの μ0 の相対標準不確かさ 2.3×10⁻¹⁰：
%%       BIPM SI Brochure 第9版（https://www.bipm.org/documents/20126/41483022/SI-Brochure-9-EN.pdf を pdftotext で読んだ）。
%%   ・エルステッドの発見（1820年。7月21日付のラテン語の4ページの報告 "Experimenta circa effectum conflictus electrici in acum magneticam"）：
%%       Smithsonian Libraries のデジタル版の説明（library.si.edu/digital-library/book/experimentacirc00orst）、Linda Hall Library
%%       "Scientist of the Day: Hans Christian Oersted"（WebSearch で照らし合わせた）。
%%   ・アンペール（1820年9月、平行な電流の引き合い・退け合い）、ビオとサバール（1820年10月）：arXiv:2102.11036（電磁誘導の歴史の論文）、
%%       EBSCO Research Starters の Ampère の項（WebSearch。日付は資料で数日ずれるので、本文は「1820年の秋」と書いた）。
%%   ・地磁気（東京）：水平分力 約30200 nT、全磁力 約46400 nT、伏角 約49°、偏角 約7°西（IGRF-11 で計算した2014年の値：
%%       北海道大学 地球惑星科学の授業資料 ep.sci.hokudai.ac.jp/~keikei/chigaku1/arch2014/No05.pdf の検索の抜粋）。数学 8-10⑪ の
%%       「日本ではおよそ 4.5×10⁻⁵ T」とも合う。〔要確認〕気象庁地磁気観測所・国土地理院・NOAA の計算サービスは作業環境から届かず、
%%       資料の本文は読めなかった。本文は丸めて「水平 約3×10⁻⁵ T、全体 約4.6×10⁻⁵ T、伏角 約50°、偏角 7〜9°西（北海道で大きい）」と書いた。
%%   ・送電線などの磁束密度 200 μT 以下（商用周波数）：電気設備に関する技術基準を定める省令 第27条の2（計画書 D12〔確〕）。
%%   ・MRI 1.5〜3 T、冷蔵庫の磁石 0.001〜0.03 T 程度：National High Magnetic Field Laboratory（MagLab）の教材
%%       "Magnets from Mini to Mighty"・"Tesla Tales"（WebSearch）。本文は「冷蔵庫の磁石の表面で 0.01 T くらい」と丸めた。
%%   ・渡り鳥の磁気の感覚：Xu ら "Magnetic sensitivity of cryptochrome 4 from a migratory songbird", Nature 594, 535–540 (2021)
%%       （オックスフォード大学のリポジトリ ora.ox.ac.uk の書誌、phys.org の解説で照らし合わせた。試験管の中のたんぱく質の実験で、
%%       鳥の目の中での確かめはまだ、と著者が述べている）。
%% 〔要確認〕：上の地磁気の値（一次資料の本文）。例題の仮の値（ホール素子の数密度 1×10²² /m³ など）は本文に「〜とします」と書いた。
\makeatletter
% 方位磁針  \sxgcompass{(中心)}{半径}{針の向き（上からの角。右回りが正）}
\newcommand\sxgcompass[3]{\begin{scope}[shift={#1}]
  \filldraw[fill=white, draw=bkInk, line width=0.6pt] (0,0) circle (#2);
  \begin{scope}[rotate={-(#3)}]
  \fill[bkRed!80!black] (0,{0.82*#2}) -- ({0.2*#2},0) -- ({-0.2*#2},0) -- cycle;
  \fill[bkGray!60] (0,{-0.82*#2}) -- ({0.2*#2},0) -- ({-0.2*#2},0) -- cycle;
  \end{scope}
  \fill[bkInk] (0,0) circle ({0.09*#2});\end{scope}}
% 紙から手前へ（⊙）・紙のおくへ（⊗）  \sxgout{(中心)}{半径}
\newcommand\sxgout[2]{\begin{scope}[shift={#1}]\draw[bkInk, line width=0.7pt, fill=white] (0,0) circle (#2); \fill[bkInk] (0,0) circle ({0.32*#2});\end{scope}}
\newcommand\sxgin[2]{\begin{scope}[shift={#1}]\draw[bkInk, line width=0.7pt, fill=white] (0,0) circle (#2); \draw[bkInk, line width=0.7pt] (45:#2) -- (225:#2) (135:#2) -- (315:#2);\end{scope}}
% 棒磁石（横向き。右が N）  \sxgmagnet{(中心)}{長さ}{太さ}
\newcommand\sxgmagnet[3]{\begin{scope}[shift={#1}]
  \filldraw[fill=bkRed!55, draw=bkInk, line width=0.6pt] (0,{-#3/2}) rectangle ({#2/2},{#3/2});
  \filldraw[fill=bkBlue!40, draw=bkInk, line width=0.6pt] ({-#2/2},{-#3/2}) rectangle (0,{#3/2});
  \node[font=\sffamily\bfseries\footnotesize] at ({#2/4},0) {N};
  \node[font=\sffamily\bfseries\footnotesize] at ({-#2/4},0) {S};\end{scope}}
% 乾電池（横向き。右が＋）  \sxgcell{(左下)}
\newcommand\sxgcell[1]{\begin{scope}[shift={#1}]
  \filldraw[fill=bkGreen!35, draw=bkInk, line width=0.6pt, rounded corners=0.6mm] (0,0) rectangle (14,6);
  \filldraw[fill=bkGray!40, draw=bkInk, line width=0.5pt] (14,1.8) rectangle (15.4,4.2);
  \node[font=\sffamily\footnotesize] at (11.6,3) {＋};\node[font=\sffamily\footnotesize] at (2.4,3) {－};\end{scope}}
% 磁場の矢印・電流の矢印・力の矢印の線の形（色だけで分けない：K-15）
\tikzset{sxg B/.style={bkBlue!80!black, line width=0.9pt, -{Stealth[length=2mm]}},
  sxg I/.style={bkOrange!85!black, line width=1.6pt, -{Stealth[length=2.6mm]}},
  sxg F/.style={bkRed!85!black, line width=1.4pt, double, double distance=0.5pt, -{Stealth[length=2.8mm]}},
  sxg lab/.style={font=\sffamily\footnotesize, text=bkInk}}
\makeatother

\phchapter{6}{電流が方位磁針を動かす：\\電流がつくる磁場}{電線に電気を流すと、なぜ方位磁針がふれる？}{\bklevelrange{1}{4}}{電流が方位磁針を動かす：電流がつくる磁場}
\bkmichishirube{4-14,8-5,8-7,8-10,8-11,力-1,力-6,電-1,電-4}{電-7,電-8,電-12,電-14}

\begin{goal}
\begin{itemize}[label={\color{bkTheme}●}, leftmargin=1.6\zw, itemsep=0pt, topsep=1pt]
\item \bklevel{1}\enspace 電流の周りに輪の形の磁場ができることを方位磁針で確かめ、その向きを右ねじで言える。
\item \bklevel{3}\enspace まっすぐな電流・円い電流・ソレノイドがつくる磁場と、磁場の中の電流・荷電粒子・コイルが受ける力を求められる。
\item \bklevel{4}\enspace ビオ・サバールの法則とアンペールの法則（積分形）から磁場を導き、$\mathrm{rot}\,\vec B=\mu_0\vec j$ と $\mathrm{div}\,\vec B=0$ を読める。
\item \bklevel{4}\enspace ホール効果・速度選別器・サイクロトロンのしくみを、ローレンツ力で説明できる。
\end{itemize}
\end{goal}

\phkoushiki{%
\phk{電流の磁界の向き・磁界から受ける力の向き}{中学理科}{ph:電-6:migineji}
\phk{円電流の中心 $B=\mu_0I/(2a)$}{高校物理}{ph:電-6:circle}
\phk{直線電流 $B=\mu_0I/(2\pi r)$}{高校物理}{ph:電-6:straight}
\phk{ソレノイド $B=\mu_0nI$（$H=nI$）}{高校物理}{ph:電-6:solenoid}
\phk{電流が受ける力 $F=IBl\sin\theta$}{高校物理}{ph:電-6:force}
\phk{平行電流の力 $F=\mu_0I_1I_2l/(2\pi r)$}{高校物理}{ph:電-6:parallel}
\phk{ローレンツ力 $f=qvB$、半径 $r=mv/(qB)$}{高校物理}{ph:電-6:lorentz}
}

\begin{hatena}
\noindent\begin{minipage}[c]{0.55\linewidth}\raggedright
机の上に方位磁針を置くと、針の赤いはし（N極）は北を向きます。その真上に、南北の向きに導線を張り、乾電池をつないで電気を流してみます。

すると、針がくるっと横を向きます。電池を外すと、針は北にもどります。導線は鉄でもないし、磁石でもありません。導線の中を流れる電気が、どうして離れた所の針を動かせるのでしょう。針は、どのくらいの角度までふれるのでしょう。
\end{minipage}\hfill
\begin{minipage}[c]{0.42\linewidth}\centering
\begin{tikzpicture}[x=1mm, y=1mm]
  \fill[bkYellow!25] (0,0) rectangle (50,40);
  \sxgcompass{(25,20)}{9}{-35}
  \draw[bkOrange!80!black, line width=1.4pt] (25,3) -- (25,37);
  \draw[sxg I] (31,10) -- (31,30);
  \node[sxg lab, anchor=west] at (32.5,26) {電流};
  \node[sxg lab, anchor=south] at (25,40.5) {北};
  \node[sxg lab, anchor=north] at (25,-0.8) {導線は針の上を通る};
  \node[sxg lab, anchor=east, text=bkRed!80!black] at (14,30) {ふれる！};
\end{tikzpicture}
\end{minipage}
\end{hatena}

\noindent 第1部（電-1〜5）では、止まった電荷がつくる電場と、導線の中を流れる電流を調べました。この章から第2部が始まります。主役は\textbf{磁場}です。磁石は、離れた鉄を引きつけ、方位磁針の向きを変えます。1820年、デンマークのエルステッドは、電流が方位磁針を動かすことを見つけました。電気と磁気は、別々のものではなかったのです。この章では、電流がどんな磁場をつくるか、そして磁場が電流や動く電荷にどんな力をおよぼすかを、実験の法則から式にしていきます。モーター・スピーカー・電磁石・MRI は、どれもこの章の式で動いています。

\begin{yattemiyou}[乾電池と導線で、エルステッドの実験をしよう]
\textbf{用意するもの}：方位磁針、単3の乾電池1本、ビニール線（\qty{50}{cm}くらい。両はしの皮をむく）、セロハンテープ\par\vspace{2pt}
\noindent\begin{minipage}[c]{0.66\linewidth}
\begin{enumerate}[label={\bkcircled[bkGreen]{\arabic*}}, leftmargin=2.4\zw, itemsep=1pt, topsep=0pt]
\item 方位磁針を机に置き、針が止まるのを待つ。導線を、針の真上に南北の向きに置き、テープでとめる。
\item 導線のはしを乾電池の＋と－に、1〜2秒だけ当てる（長く当てると電池と導線が熱くなる）。針はどちらへふれるか。
\item 電池の向きを反対にして、②をする。
\item 導線を方位磁針の下にくぐらせて、②をする。
\end{enumerate}
\end{minipage}\hfill
\begin{minipage}[c]{0.32\linewidth}\centering
\begin{tikzpicture}[x=0.85mm, y=0.85mm]
  \sxgcompass{(20,16)}{8}{0}
  \draw[bkOrange!80!black, line width=1.2pt] (20,2) -- (20,30) -- (36,30) -- (36,22);
  \draw[bkOrange!80!black, line width=1.2pt] (20,2) -- (36,2) -- (36,9.6);
  \begin{scope}[rotate around={90:(36,16)}]\sxgcell{(29,13)}\end{scope}
  \node[sxg lab, anchor=south] at (20,30.6) {北};
\end{tikzpicture}
\end{minipage}
\par\vspace{2pt}\noindent\textbf{注意}：家のコンセントは、ぜったいに使わない。電池は1本だけ、短い時間だけ当てる。
\end{yattemiyou}

\begin{bkcenter}
\begin{bktable}{colspec={Q[c,m,wd=12mm] Q[c,m,wd=30mm] Q[c,m,wd=30mm] Q[c,m,wd=36mm]}, rowsep=1pt, row{1}={bg=bkTheme!10}}
手順 & 電流の向き & 導線の場所 & 針のN極がふれた向き \\
② & 南から北へ & 針の上 & 西へ約30° \\
③ & 北から南へ & 針の上 & 東へ約30° \\
④ & 南から北へ & 針の下 & 東へ約30° \\
\end{bktable}
\bktabcaption{やってみようで、ソラさんが調べた結果（ふれの角は目で見た大きさ）}\label{ph:電-6:tab-yattemiyou}
\end{bkcenter}

\noindent やってみようの答え：表\ref{ph:電-6:tab-yattemiyou}のように、南から北へ電流を流すと、上から見て、針のN極は西（左）へふれます。電池の向きを反対にすると、東（右）へふれます。導線を針の下にくぐらせると、同じ向きの電流でも、ふれる向きが反対になります。乾電池1本の電流は、導線の長さや電池の状態で変わりますが、1〜数 A くらいです。電流をそのままにして導線を針から遠ざけると、ふれは小さくなります。この「上と下で向きが反対」が、大事な手がかりです。導線の周りの磁場は、導線を中心に\textbf{ぐるっと回る輪の形}をしているのです（②）。どのくらいふれるかは、⑤で数を使って求めます。

\begin{omoidasou}
\bkrecall{4-14}{外積 $\vec a\times\vec b$ は、$\vec a$ にも $\vec b$ にも垂直で、$\vec a$ から $\vec b$ へ回す右ねじが進む向き。大きさ $ab\sin\theta$。紙の手前へ向かう矢印は $\odot$、おくへ向かう矢印は $\otimes$（4-14④）。}
\bkrecall{4-14}{磁場 $\vec B$ の中を速度 $\vec v$ で動く電荷 $q$ は、力 $\vec F=q\vec v\times\vec B$（ローレンツ力）を受ける。導線では $\vec F=I\vec L\times\vec B$。\qty{3}{A}・\qty{0.2}{m}・\qty{0.5}{T}で\qty{0.3}{N}（4-14⑧）。}
\bkrecall{8-10}{まっすぐな電流の周りの磁場は $B=\mu_0I/(2\pi r)$。\qty{10}{A}から\qty{5}{cm}で $4\times10^{-5}$\,T（8-10⑪）。}
\bkrecall{8-11}{アンペールの法則 $\oint_C\vec B\cdot\dd\vec r=\mu_0I_{\text{中}}$。対称な場では、楽な道を選ぶ（8-11⑬）。}
\bkrecall{電-1}{電場は、試しの電荷 $q$ が受ける力を $q$ でわったもの（$\vec E=\vec F/q$）。点電荷の電場は距離の2乗に反比例（電-1 ④）。}
\bkrecall{電-4}{電流は1秒に通る電荷の量。導線の中では、電子が電流と反対向きに流れる（電-4 ①②）。}
\bkrecall{力-1}{速さ $v$、半径 $r$ の等速円運動の加速度は、中心の向きに $v^2/r$（力-1 ⑤）。}
\end{omoidasou}

\noindent この章の記号：磁場（磁束密度）は $\vec B$、導線の長さは $l$ と書きます（数学 4-14 の $\vec L$ と同じもの）。トルクは $T$ と書きます（第1巻の $\vec N$ と同じもの。電-12 で詳しく使います）。

%% ======================== 物理の道具 ========================
\bksection{doguu}{物理の道具}

\bksubsection{1}{磁石と磁場：磁力線で見る}
\label{ph:電-6:sec1}
磁石は、鉄を引きつけます。2つの磁石を近づけると、N極とS極は引き合い、N極どうし・S極どうしは退け合います。小学校の理科で調べたことです。磁石の力は、間に紙やガラスをはさんでも、離れていても届きます。電-1の電気の力と同じように、磁石の周りの空間そのものが、「ここに方位磁針を置くと、この向きに回される」という性質をもっていると考えます。この空間の性質を\textbf{磁場}といいます。\bkidx{じば}{磁場}

方位磁針は、それ自身が小さな磁石です。磁場の中に置くと、針のN極が向く向きが、その場所の\textbf{磁場の向き}です。いろいろな場所に方位磁針を置いて、針の向きをつないでいくと、線ができます。これを\textbf{磁力線}といいます（8-5④の力線）。\bkidx{じりょくせん}{磁力線}

\begin{bkfigure}
\begin{tikzpicture}[x=1mm, y=1mm]
  \path[use as bounding box] (-34,-32) rectangle (92,32);
  \foreach \kk/\hh in {6/8,12/16,18/26,24/38} {
    \draw[bkBlue!70!black, line width=0.6pt] (11,1.5) .. controls ({11+\kk},{1.5+\hh}) and ({-11-\kk},{1.5+\hh}) .. (-11,1.5);
    \draw[bkBlue!70!black, line width=0.6pt] (11,-1.5) .. controls ({11+\kk},{-1.5-\hh}) and ({-11-\kk},{-1.5-\hh}) .. (-11,-1.5);
    \draw[bkBlue!80!black, line width=0.8pt, -{Stealth[length=2mm]}] (1.2,{1.5+0.75*\hh}) -- (-1.2,{1.5+0.75*\hh});
    \draw[bkBlue!80!black, line width=0.8pt, -{Stealth[length=2mm]}] (1.2,{-1.5-0.75*\hh}) -- (-1.2,{-1.5-0.75*\hh});}
  \draw[bkBlue!70!black, line width=0.6pt] (11,0) -- (21,0) (27,0) -- (32,0) (-32,0) -- (-27,0) (-21,0) -- (-11,0);
  \draw[bkBlue!80!black, line width=0.8pt, -{Stealth[length=2mm]}] (-32,0) -- (-28.5,0);
  \draw[bkBlue!80!black, line width=0.8pt, -{Stealth[length=2mm]}] (28,0) -- (32,0);
  \sxgmagnet{(0,0)}{22}{6}
  \sxgcompass{(24,0)}{3}{90}
  \sxgcompass{(-24,0)}{3}{90}
  \sxgcompass{(0,17.3)}{2.5}{-90}
  \node[anchor=west, font=\small, text width=54mm, align=left] at (38,8) {%
  磁力線は、磁石の外では\textbf{N極から出て、S極に入る}。\par\vspace{2pt}
  方位磁針のN極（赤いほう）は、磁力線の向きを向く。\par\vspace{2pt}
  磁力線が混んでいる所（極の近く）ほど、磁場が強い。\par\vspace{2pt}
  磁力線は、とちゅうで切れたり、交わったりしない。};
\end{tikzpicture}
\bkfigcaption{棒磁石の周りの磁力線と、方位磁針の向き。砂鉄をまくと、この形の模様が見える。}\label{ph:電-6:fig-bar}
\end{bkfigure}

\noindent 磁力線は、電-1の電気力線とよく似ています。違うのは、磁石を2つに切っても、N極だけ・S極だけの磁石は取り出せないことです。切った所に新しいS極とN極ができ、2つの磁石になります。磁力線は、磁石の中ではS極からN極へ進み、\textbf{いつも輪になってつながって}います（8-9 の $\mathrm{div}\,\vec B=0$。⑦で式にします）。

\begin{phmijika}[地球は大きな磁石]
方位磁針のN極が北を向くのは、地球そのものが大きな磁石だからです。北極の近くに地球の磁石のS極が、南極の近くにN極があります（N極を引きつけるのはS極）。地球の磁場は、⑫で詳しく見ます。
\end{phmijika}

\bksubsection{2}{エルステッドの実験と右ねじの法則}
\label{ph:電-6:sec2}
1820年、デンマークのエルステッドは、電流を流した導線の近くで方位磁針がふれることに気づき、その年の7月に、ラテン語の4ページの報告にまとめて各地の研究者に送りました。やってみようと同じ実験です。導線の真上と真下で、ふれる向きは反対でした。導線を囲むようにいくつも方位磁針を並べると、針は\textbf{導線を中心とする円の接線の向き}を向きます。

\begin{bkfigure}
\begin{tikzpicture}[x=1mm, y=1mm]
  % 左：横から見た図
  \foreach \ry in {6,16} {\draw[bkBlue!70!black, line width=0.7pt] (36,\ry) arc (0:180:16 and 4.5);}
  \draw[bkOrange!80!black, line width=2pt] (20,-6) -- (20,30);
  \foreach \ry in {6,16} {\draw[bkBlue!70!black, line width=0.7pt] (4,\ry) arc (180:360:16 and 4.5);}
  \draw[sxg I] (24,22) -- (24,30);
  \node[sxg lab, anchor=west] at (25,27) {電流 $I$};
  \draw[sxg B] (13,1.65) -- (16,1.55);
  \draw[sxg B] (27,10.35) -- (24,10.45);
  \draw[sxg B] (13,11.65) -- (16,11.55);
  \node[sxg lab, anchor=north] at (20,-7) {横から見る};
  \node[sxg lab, anchor=west, text=bkBlue!70!black] at (36.5,16) {$\vec B$};
  % まん中：上から見た図（電流は手前へ）
  \begin{scope}[shift={(63,12)}]
  \foreach \rr in {6,11,16} {\draw[bkBlue!70!black, line width=0.6pt] (0,0) circle (\rr);}
  \foreach \aa in {0,90,180,270} {\draw[sxg B] (\aa:11) ++(\aa-90:1.5) -- ++(\aa+90:3);}
  \sxgout{(0,0)}{2.4}
  \sxgcompass{(0,-16)}{2.8}{90}
  \sxgcompass{(16,0)}{2.8}{0}
  \sxgcompass{(0,16)}{2.8}{-90}
  \sxgcompass{(-16,0)}{2.8}{180}
  \node[sxg lab, anchor=north] at (0,-20) {上から見る（電流は手前へ $\odot$）};
  \end{scope}
  % 右：右ねじ
  \begin{scope}[shift={(103,4)}]
  \fill[bkOrange!25] (-3,0) rectangle (3,22);
  \draw[bkInk, line width=0.6pt] (-3,0) rectangle (3,22);
  \foreach \yy in {2,5,...,20} {\draw[bkInk!60, line width=0.4pt] (-3,\yy) -- (3,\yy+1.5);}
  \draw[sxg I] (0,23) -- (0,30);
  \node[sxg lab, anchor=south] at (0,30.2) {進む向き＝電流};
  \draw[bkBlue!80!black, line width=1pt, -{Stealth[length=2.2mm]}] (-7,9) arc (200:520:7 and 2.4);
  \node[sxg lab, anchor=north, align=center] at (0,-1) {回す向き\\＝磁場};
  \end{scope}
\end{tikzpicture}
\bkfigcaption{まっすぐな電流の周りの磁場。磁力線は導線を囲む円で、向きは「電流の向きに進む右ねじを回す向き」。上から見て、電流が手前へ（$\odot$）向かうとき、磁場は左回り。}\label{ph:電-6:fig-oersted}
\end{bkfigure}

\noindent\label{ph:電-6:migineji}この向きの決まりを\textbf{右ねじの法則}といいます。\bkidx{みぎねじのほうそく}{右ねじの法則}「電流の向きに右ねじを進めるとき、ねじを回す向きが磁場の向き」です。右手の親指を電流の向きに立てて、ほかの4本の指を曲げると、指の先が磁場の向きを指します。4-14④で、「回る向きを、軸の向きの矢印で表す」ことを学びました。右ねじの法則は、その反対の読み方です。軸の矢印（電流）が先にあって、周りを回る向き（磁場）が決まります。

やってみようで、南から北へ電流を流すと、導線の真下にある針の所では、磁場は西を向きます。北に立って南を見ると、電流は手前へ向かうので、図\ref{ph:電-6:fig-oersted}のまん中と同じ向きです。導線の真下の点では磁場は右向きで、南を見ているときの右は西です。だから針のN極は西へふれました。導線を針の下にくぐらせると、針は導線の上にあるので、磁場は東を向き、反対にふれます。

\begin{phmijika}[電気のコードのそばの方位磁針]
家の電気のコードは、行きと帰りの2本の線が並んでいます。2本の電流は反対向きなので、それぞれの磁場は、少し離れるとほとんど打ち消し合います（例題12）。コードのそばに方位磁針を置いても、針がほとんど動かないのはこのためです。
\end{phmijika}

\bksubsection{3}{磁場の強さを力で決める：ローレンツ力と単位 T}
\label{ph:電-6:sec3}
電-1では、試しの電荷 $q$ が受ける力から、電場を $\vec E=\vec F/q$ と決めました。磁場の強さも、力で決めます。ただし、磁場は\textbf{止まった電荷には力をおよぼしません}。動いている電荷だけが力を受け、その力は次の3つの性質をもちます（たくさんの実験からわかったこと）。
\begin{enumerate}[label={\bkcircled{\arabic*}}, leftmargin=2.4\zw, itemsep=1pt, topsep=2pt]
\item 力の大きさは、電荷 $q$ と速さ $v$ に比例する。
\item 速度が磁場の向きと平行なら、力は0。直角のとき、いちばん大きい。
\item 力の向きは、速度にも磁場の向きにも垂直。
\end{enumerate}
この3つをまとめて書けるのが、4-14⑧の外積の式です。
\pheq{\vec F=q\,\vec v\times\vec B}{N ＝ C × m/s × T}
\label{ph:電-6:lorentzdef}この力を\textbf{ローレンツ力}といいます。\bkidx{ろーれんつりょく}{ローレンツ力}大きさは $F=qvB\sin\theta$（$\theta$ は $\vec v$ と $\vec B$ の間の角）です。この式で、磁場 $\vec B$ の大きさと向きが決まります。単位は\textbf{テスラ}〔T〕で、\qty{1}{C}の電荷が\qty{1}{m/s}で磁場に直角に動くとき\qty{1}{N}の力を受ける磁場が\qty{1}{T}です。$\mathrm{C}\cdot\mathrm{m/s}=\mathrm{A}\cdot\mathrm{m}$ なので、$\qty{1}{T}=\qty{1}{N/(A.m)}$ です。\bkidx{てすら}{テスラ（T）}

\begin{bkfigure}
\begin{tikzpicture}[x=1mm, y=1mm]
  \draw[bkGray, line width=0.5pt, -{Stealth[length=2mm]}] (0,0) -- (40,0);
  \draw[bkGray, line width=0.5pt, -{Stealth[length=2mm]}] (0,0) -- (0,32);
  \draw[sxg B, line width=1.4pt] (0,0) -- (30,0);
  \node[sxg lab, anchor=north] at (28,-1) {磁場 $\vec B$};
  \draw[sxg I] (0,0) -- (-14,-11);
  \node[sxg lab, anchor=north] at (-14,-12) {速度 $\vec v$};
  \draw[sxg F] (0,0) -- (0,26);
  \node[sxg lab, anchor=west] at (1.5,24) {力 $\vec F=q\vec v\times\vec B$};
  \fill[bkRed] (0,0) circle (1.4);
  \node[sxg lab, anchor=north west] at (1,-1) {$+q$};
  \node[anchor=west, font=\small, text width=60mm, align=left] at (44,12) {%
  $\vec v$ から $\vec B$ へ、小さいほうの角で回す右ねじが進む向きが、$q>0$ のときの力の向き。\par\vspace{2pt}
  中学で学んだフレミングの左手の法則（中指＝電流、人さし指＝磁場、親指＝力）と、いつも同じ向き（4-14⑧）。\par\vspace{2pt}
  $q<0$（電子）なら、力は反対向き。};
\end{tikzpicture}
\bkfigcaption{ローレンツ力の向き。力はいつも速度に垂直なので、仕事をしない（速さを変えず、向きだけを変える）。}\label{ph:電-6:fig-lorentz}
\end{bkfigure}

\begin{phyomi}[$\vec F=q\vec v\times\vec B$]
「磁場は、動く電荷を、横へ押す」。電場の力 $q\vec E$ は、電荷が止まっていてもはたらき、電場の向きに押します。磁場の力は、動いていないとはたらかず、押す向きは、動く向きにも磁場にも\textbf{横向き}です。力がいつも速度に垂直なので、ローレンツ力は電荷に仕事をしません（力-4 の仕事は、力と動いた向きの内積）。磁場は電荷の速さを変えず、進む向きだけを曲げます（⑩）。
\end{phyomi}

\phboxhead{真空の透磁率 $\mu_0$}
電-1では、クーロンの法則の比例の定数を $1/(4\pi\varepsilon_0)$ と書き、$\varepsilon_0$ を真空の誘電率とよびました。電流がつくる磁場の式には、同じように\textbf{真空の透磁率} $\mu_0$（ミューゼロ）が入ります（④〜⑥）。\bkidx{しんくうのとうじりつ}{真空の透磁率}
\pheq{\mu_0=4\pi\times10^{-7}\ \mathrm{N/A^2}\approx1.26\times10^{-6}\ \mathrm{N/A^2}}{N/A$^2$ ＝ T·m/A}
2019年5月19日までは、アンペアを「平行な2本の導線の間の力」で決めていたので（⑨）、$\mu_0$ はちょうど $4\pi\times10^{-7}$ と決まった数でした。2019年5月20日から、アンペアは電気素量 $e=1.602\,176\,634\times10^{-19}$\,C を決まった数にして決めるようになり、$\mu_0$ は\textbf{実験で測って決める数}になりました。今の値（2022年に国際的にまとめられた値）は $1.256\,637\,061\,27\times10^{-6}$\,N/A$^2$ で、$4\pi\times10^{-7}$ との違いは1億分の1のさらに100分の1くらい（相対 $1\times10^{-10}$ ほど）です。この本では $\mu_0=4\pi\times10^{-7}$ N/A$^2$ を使います。

\phboxhead{高校の教科書の $H$ と、この本の $B$}
\label{ph:電-6:BH}高校の物理の教科書は、電流がつくる磁場を、まず「磁場（磁界）の強さ」$H$〔A/m〕で書き（たとえば $H=I/(2\pi r)$）、それに真空の透磁率をかけた
\pheq{B=\mu_0H}{T ＝ N/A$^2$ × A/m}
を「磁束密度」とよびます。この本は、数学の本（8-10・8-14）と同じく、力を決める $\vec B$ のほうを基本の場とし、「磁場（磁束密度）」とよびます。真空（と空気）の中では、$B$ と $H$ は $\mu_0$ 倍の違いしかないので、どちらで考えても同じです。鉄のような物の中に入ると、2つは別の役わりをもつようになります。$H$ は電-7 で「磁場の強さ」として、もう一度出てきます。\bkidx{じそくみつど}{磁束密度}

\begin{phketa}[磁場の強さの目安]
\phketaline{地球の磁場（日本）}{約 $4.6\times10^{-5}$\,T（$\qty{46}{\micro T}$）}
\phketaline{\qty{10}{A}の電流から\qty{5}{cm}}{$4\times10^{-5}$\,T（地球と同じくらい：⑤）}
\phketaline{冷蔵庫に貼る磁石の表面}{\qty{0.01}{T}くらい}
\phketaline{病院の MRI}{1.5〜\qty{3}{T}（地球の磁場の3万〜6万倍）}
\phboxhead{答えが桁で合っているか}
身の回りの電流（数 A）が数 cm の所につくる磁場は、地球の磁場と同じくらいの $10^{-5}$\,T の桁です。電流の磁場の計算で $1$\,T の桁が出たら、何かをまちがえています（ソレノイドや鉄心を使わない限り）。
\end{phketa}

\begin{senmonsho}[電気工学の専門書の書き方]
\begin{bkhonbun}
無限長直線電流 $I$ から距離 $r$ の点の\textbf{磁界の強さ}は $H=I/2\pi r$〔A/m〕であり，真空中の\textbf{磁束密度}は $B=\mu_0H$〔T〕である．ここで $\mu_0=4\pi\times10^{-7}$〔H/m〕は真空の透磁率である．電界 $\symbfit E$，磁束密度 $\symbfit B$ 中を速度 $\symbfit v$ で運動する電荷 $q$ に働く力は $\symbfit F=q(\symbfit E+\symbfit v\times\symbfit B)$ である．
\end{bkhonbun}
\bkyaku{電気工学の本は、電場を「電界」、磁場（$H$）を「磁界」とよぶ（場＝界）。「磁界の強さ」$H$ は高校の $H$ と同じ。この本の「磁場（磁束密度）」$B$ は「磁束密度」。$\mu_0$ の単位 H/m（ヘンリー毎メートル）は、この本の N/A$^2$ と同じもの（H はインダクタンスの単位：電-8）。太字はベクトル。最後の式は、電場の力とローレンツ力をまとめた形。古い本や物理の論文では、磁場をガウス（G）で書くことがある。$\qty{1}{T}=10^4$\,G。}
\end{senmonsho}

\bksubsection{4}{ビオ・サバールの法則：円い電流の中心}
\label{ph:電-6:sec4}
エルステッドの発見が伝わると、1820年の秋のうちに、フランスのビオとサバールは、導線からの距離を変えて方位磁針のふれをはかり、電流のつくる磁場の式を見つけました。その考えは、「導線を短い部分に分け、1つ1つの部分がつくる磁場をたし合わせる」というものです（電-1 ⑦で、広がった電荷の電場を、小さな電荷の電場のたし合わせで求めたのと同じ考え）。

\begin{bkfigure}
\begin{tikzpicture}[x=1mm, y=1mm]
  \draw[bkOrange!80!black, line width=1.4pt] (0,0) .. controls (12,6) and (24,18) .. (30,34);
  \draw[sxg I] (12.5,7.2) -- (18.5,12.2);
  \node[sxg lab, anchor=north west] at (17,9.5) {$I\,\dd\vec l$};
  \fill[bkInk] (12.5,7.2) circle (0.8);
  \fill[bkRed] (62,10) circle (1.2);
  \node[sxg lab, anchor=west] at (63,9) {点P};
  \draw[bkGray, line width=0.6pt, -{Stealth[length=2mm]}] (12.5,7.2) -- (61,9.9);
  \node[sxg lab, anchor=south] at (38,9.2) {距離 $r$、向き $\vec e_r$};
  \draw[bkInk, line width=0.5pt] (16.5,10.6) arc (40:5:5);
  \node[sxg lab, anchor=west] at (18.8,12.4) {$\theta$};
  \sxgin{(62,3)}{2.2}
  \node[sxg lab, anchor=east] at (59,2) {$\dd\vec B$（おくへ）};
  \node[anchor=west, font=\small, text width=46mm, align=left] at (84,12) {%
  向き：$\dd\vec l\times\vec e_r$（$\dd\vec l$ と $\vec e_r$ の両方に垂直）。\par\vspace{2pt}
  大きさ：電流と長さに比例し、距離の2乗に反比例、$\sin\theta$ に比例。};
\end{tikzpicture}
\bkfigcaption{ビオ・サバールの法則。導線の短い部分 $\dd\vec l$ を流れる電流が、点Pにつくる磁場 $\dd\vec B$。}\label{ph:電-6:fig-bs}
\end{bkfigure}

\noindent 導線の短い部分（長さ $\dd l$、電流の向きのベクトル $\dd\vec l$）を流れる電流 $I$ が、距離 $r$ の点Pにつくる磁場は、
\pheq{\dd\vec B=\frac{\mu_0}{4\pi}\,\frac{I\,\dd\vec l\times\vec e_r}{r^2}\qquad\Bigl(\text{大きさ}\ \dd B=\frac{\mu_0}{4\pi}\,\frac{I\,\dd l\,\sin\theta}{r^2}\Bigr)}{T ＝ $\dfrac{\text{N}}{\text{A}^2}$×$\dfrac{\text{A·m}}{\text{m}^2}$}
\label{ph:電-6:biot}です（$\vec e_r$ は、その部分から点Pへ向かう長さ1のベクトル）。これを\textbf{ビオ・サバールの法則}といいます。\bkidx{びお・さばーるのほうそく}{ビオ・サバールの法則}導線全体の磁場は、これを導線にそって積分（たし合わせ）すると求められます。

\begin{phyomi}[$\dd B=\dfrac{\mu_0}{4\pi}\dfrac{I\,\dd l\sin\theta}{r^2}$]
「電流の小さな部分は、距離の2乗に反比例する磁場を、横向きにつくる」。距離の2乗に反比例するのは、クーロンの法則と同じです。違うのは、向きが $\dd\vec l\times\vec e_r$ で、電流の向きにも、点Pの向きにも垂直なことです。電流の向きの延長線の上（$\theta=0$）には、磁場をつくりません。
\end{phyomi}

\phboxhead{円い電流の中心}
\label{ph:電-6:circle}半径 $a$ の円の形の導線に電流 $I$ を流します。円の中心では、導線のどの部分も、中心までの距離は $a$ で、$\dd\vec l$ と中心へ向かう向き $\vec e_r$ は直角（$\sin\theta=1$）です。しかも、どの部分がつくる磁場も、円の面に垂直な同じ向きです（右ねじの向き）。だから、大きさをそのままたせばよく、円周の長さ $2\pi a$ の分を足すと、
\[B=\frac{\mu_0}{4\pi}\,\frac{I\times2\pi a}{a^2}\]
\pheq{B=\frac{\mu_0I}{2a}}{T ＝ N/A$^2$ × A÷m}
です。$N$ 回巻いたコイルなら、$N$ 倍の $B=\mu_0NI/(2a)$ です。高校の教科書の $H=I/(2a)$ に $\mu_0$ をかけたものです。

\begin{bkfigure}
\begin{tikzpicture}[x=1mm, y=1mm]
  \draw[bkOrange!80!black, line width=1.6pt] (0,0) ellipse (22 and 7);
  \draw[sxg I] (-4,-7) -- (4,-7);
  \node[sxg lab, anchor=north] at (0,-8) {電流 $I$};
  \draw[bkBlue!70!black, line width=1.4pt, -{Stealth[length=3mm]}] (0,-2) -- (0,24);
  \node[sxg lab, anchor=west] at (1.5,22) {$\vec B$（中心で $\mu_0I/(2a)$）};
  \draw[bkGray, line width=0.5pt, dashed] (0,0) -- (22,0);
  \node[sxg lab, anchor=south] at (13,0.4) {半径 $a$};
  \foreach \xx in {-22,22} {\draw[bkBlue!60!black, line width=0.6pt] (\xx,0) ellipse (3.6 and 8.5);}
  \node[anchor=west, font=\small, text width=56mm, align=left] at (36,4) {%
  上から見て左回りに流れる電流は、上向きの磁場をつくる（右手の4本の指を電流の向きに曲げると、親指が磁場の向き）。\par\vspace{2pt}
  コイルの上の面がN極、下の面がS極の、平たい磁石のようになる。};
\end{tikzpicture}
\bkfigcaption{円い電流がつくる磁場。導線の近くでは、導線を囲む小さな輪の磁力線。中心ではどの部分の磁場も同じ向きにそろう。}\label{ph:電-6:fig-circle}
\end{bkfigure}

\begin{phbaisuru}{半径 $a$}
中心の磁場は半分になります。導線の長さは2倍になりますが、どの部分も2倍遠く（磁場は $\tfrac14$）なるので、$2\times\tfrac14=\tfrac12$ です。小さなコイルほど、中心の磁場は強くなります。
\end{phbaisuru}

\bksubsection{5}{まっすぐな電流：はてな？の答え}
\label{ph:電-6:sec5}
とても長いまっすぐな導線（$z$ 軸）に電流 $I$ を流し、導線から距離 $r$ の点Pの磁場を求めます。導線の $z$ の所にある短い部分 $\dd z$ から点Pまでの距離を $s$ とすると、$s^2=r^2+z^2$、$\sin\theta=r/s$ です（図\ref{ph:電-6:fig-straight}）。どの部分がつくる磁場も、点Pで同じ向き（導線を回る向き）なので、大きさをたします。
\[B=\int_{-\infty}^{\infty}\frac{\mu_0}{4\pi}\,\frac{I\,\dd z}{s^2}\cdot\frac{r}{s}=\frac{\mu_0Ir}{4\pi}\int_{-\infty}^{\infty}\frac{\dd z}{(r^2+z^2)^{3/2}}\]
この積分は、$z=r\tan\alpha$ と角でおきかえると（5-16⑥）、
\[\dd z=\frac{r}{\cos^2\alpha}\,\dd\alpha,\qquad(r^2+z^2)^{3/2}=\frac{r^3}{\cos^3\alpha}\]
なので、
\[\int_{-\infty}^{\infty}\frac{\dd z}{(r^2+z^2)^{3/2}}=\int_{-\pi/2}^{\pi/2}\frac{\cos\alpha}{r^2}\,\dd\alpha=\frac{1}{r^2}\Bigl[\sin\alpha\Bigr]_{-\pi/2}^{\pi/2}=\frac{2}{r^2}\]
です。だから、
\pheq{B=\frac{\mu_0I}{2\pi r}}{T ＝ N/A$^2$ × A÷m}
\label{ph:電-6:straight}です。8-10⑪で「物理の実験からわかっている」とした式が、ビオ・サバールの法則から導けました。高校の教科書の $H=I/(2\pi r)$ に $\mu_0$ をかけたものです。\bkidx{ちょくせんでんりゅう}{直線電流の磁場}

\begin{bkfigure}
\begin{tikzpicture}[x=1mm, y=1mm]
  \draw[bkOrange!80!black, line width=1.6pt] (0,-14) -- (0,34);
  \draw[sxg I] (-3.5,-10) -- (-3.5,2);
  \node[sxg lab, anchor=east] at (-5.5,-4) {$I$};
  \fill[bkInk] (0,18) circle (0.9);
  \draw[bkOrange!90!black, line width=3pt] (0,16.5) -- (0,19.5);
  \node[sxg lab, anchor=east] at (-1.5,18) {$\dd z$};
  \draw[bkGray, line width=0.5pt, dashed] (0,0) -- (40,0);
  \fill[bkRed] (40,0) circle (1.2);
  \node[sxg lab, anchor=west] at (41.5,0) {点P};
  \draw[bkGray, line width=0.6pt] (0,18) -- (40,0);
  \node[sxg lab, anchor=south west] at (21,9.5) {$s=\sqrt{r^2+z^2}$};
  \node[sxg lab, anchor=north] at (20,-0.6) {$r$};
  \draw[bkInk, line width=0.5pt, {Stealth[length=1.6mm]}-{Stealth[length=1.6mm]}] (-8,0.3) -- (-8,17.7);
  \node[sxg lab, anchor=east] at (-9,9) {$z$};
  \draw[bkInk, line width=0.5pt] (0,13) arc (-90:-24:5);
  \node[sxg lab, anchor=north west] at (3.6,13.4) {$\theta$};
  \draw[bkInk, line width=0.5pt] (34,0) arc (180:155:6);
  \node[sxg lab, anchor=south east] at (34,0.3) {$\alpha$};
  \sxgin{(40,7)}{2}
  \node[sxg lab, anchor=west] at (43,7) {$\dd\vec B$};
  \node[anchor=west, font=\small, text width=62mm, align=left] at (56,12) {%
  上の部分も下の部分も、点Pに同じ向き（ここでは紙のおくへ）の磁場をつくる。\par\vspace{2pt}
  $\alpha$ は、点Pから見た $\dd z$ の高さの角。$z=r\tan\alpha$。導線のはしからはしまでで、$\alpha$ は $-90^\circ$ から $90^\circ$ まで動く。};
\end{tikzpicture}
\bkfigcaption{まっすぐな電流の磁場を、ビオ・サバールの法則で求める。}\label{ph:電-6:fig-straight}
\end{bkfigure}

\begin{phyomi}[$B=\dfrac{\mu_0I}{2\pi r}$]
「磁場は電流に比例し、導線からの距離に反比例する。磁力線は導線を囲む円」。$2\pi r$ は、半径 $r$ の円の周りの長さです（⑥で、この理由がわかります）。\qty{1}{A}の電流から\qty{1}{m}の所では $2\times10^{-7}$\,T で、地球の磁場の約200分の1です。
\end{phyomi}

\begin{phbaisuru}{距離 $r$}
磁場は半分になります。点電荷の電場（電-1 ④）は、距離が2倍で $\tfrac14$ でした。まっすぐな電流の磁場は、それよりゆっくり弱まります。下のグラフは、\qty{10}{A}の電流の磁場です。
\par\smallskip\noindent\begin{tikzpicture}
\begin{axis}[bkplot, width=96mm, height=42mm, xmin=0, xmax=0.23, ymin=0, ymax=60, xtick={0.05,0.10,0.15,0.20}, xticklabel style={/pgf/number format/fixed}, ytick={10,20,40}, xlabel={距離 $r$（m）}, ylabel={$B$（$\mu$T）}]
  \addplot[bkcurve, domain=0.035:0.22, samples=80] {2/x};
  \addplot[bkOrange, line width=1.3pt, dash pattern=on 4pt off 2pt, domain=0.045:0.22, samples=80] {40*(0.05/x)^2};
  \fill[bkRed] (axis cs:0.05,40) circle (1.8pt) (axis cs:0.10,20) circle (1.8pt) (axis cs:0.20,10) circle (1.8pt);
  \node[font=\sffamily\footnotesize, anchor=south west, text=bkRed!75!black] at (axis cs:0.105,22) {距離2倍で半分};
  \node[font=\sffamily\footnotesize, anchor=west, fill=white, inner sep=1pt] at (axis cs:0.115,50) {破線：$1/r^2$ で弱まるとしたら};
\end{axis}
\end{tikzpicture}
\end{phbaisuru}

\phboxhead{\bkhatenaref の答え}
\label{ph:電-6:hatena}導線の電流は、その周りに輪の形の磁場をつくります（②）。方位磁針は、地球の磁場と、電流の磁場を\textbf{ベクトルとしてたした}磁場の向きを向きます。日本の地面の近くでは、地球の磁場の水平の部分は、北向きに約 $3.0\times10^{-5}$\,T です（⑫）。\qty{10}{A}の電流が、針から\qty{5}{cm}の所を南北に流れていると、針の所の電流の磁場は
\[B=\frac{\mu_0I}{2\pi r}=\frac{4\pi\times10^{-7}\times10}{2\pi\times0.05}=4.0\times10^{-5}\ \text{T}\]
で、東西の向きです。北向き $3.0$、東西向き $4.0$（単位 $10^{-5}$\,T）のベクトルをたすと、針は北から $\tan\theta=4.0/3.0$ の角、約53°ふれます（例題1）。電流の磁場が地球の磁場と同じくらいの強さだから、針は大きく向きを変えるのです。

\begin{bkfigure}
\begin{tikzpicture}[x=0.8mm, y=0.8mm]
  \draw[bkBlue!80!black, line width=1.4pt, -{Stealth[length=3mm]}] (0,0) -- (0,30);
  \node[sxg lab, anchor=west] at (1,20) {地球 $3.0$};
  \draw[bkOrange!85!black, line width=1.4pt, dash pattern=on 4pt off 2pt, -{Stealth[length=3mm]}] (0,30) -- (-40,30);
  \node[sxg lab, anchor=south] at (-20,31) {電流の磁場 $4.0$（西）};
  \draw[bkRed!80!black, line width=1.8pt, -{Stealth[length=3mm]}] (0,0) -- (-40,30);
  \node[sxg lab, anchor=north east, align=right] at (-21,14) {合わせた\\磁場 $5.0$};
  \draw[bkInk, line width=0.5pt] (0,10) arc (90:143.13:10);
  \node[sxg lab, anchor=south east] at (-2,10) {53°};
  \sxgcompass{(0,0)}{5}{-53}
  \node[sxg lab] at (-20,-8) {（単位 $10^{-5}$\,T。上が北）};
\end{tikzpicture}\hspace{4mm}
\begin{tikzpicture}
\begin{axis}[bkplot, width=58mm, height=44mm, xmin=0, xmax=52, ymin=0, ymax=92, xtick={10,20,30,40,50}, ytick={30,60,90}, xlabel={電流（A）}, ylabel={ふれ（度）}]
  \addplot[bkcurve, domain=0:50, samples=80] {atan(0.4*x/3)};
  \fill[bkRed] (axis cs:10,53.13) circle (1.8pt);
  \node[font=\sffamily\footnotesize, anchor=north west, text=bkRed!75!black] at (axis cs:11,50) {\qty{10}{A}で53°};
\end{axis}
\end{tikzpicture}
\bkfigcaption{はてな？の答え。左：地球の磁場と電流の磁場をたす。右：針から\qty{5}{cm}の所の電流と、ふれの角。電流を大きくしても90°にはならない（地球の磁場は消えないから）。}\label{ph:電-6:fig-hatena}
\end{bkfigure}

\begin{phmodel}[導線の真下の方位磁針]
\phignore{電池や、つないだ導線のほかの部分の電流がつくる磁場。針が小さな磁石として周りにおよぼす影響。}
\phassume{導線はまっすぐで、針までの距離より十分長い。針の所の磁場は、地球の磁場の水平の部分と電流の磁場の和。}
\phlimit{導線が短く曲がっているとき（④のビオ・サバールで積分しなおす）。近くに鉄の机・スピーカーなどがあるとき。}
\end{phmodel}

\bksubsection{6}{アンペールの法則：道を選んで磁場を求める}
\label{ph:電-6:sec6}
⑤では、ビオ・サバールの法則を積分して、まっすぐな電流の磁場を求めました。磁場の形に対称性があるときは、もっと楽な方法があります。8-11⑬で学んだ\textbf{アンペールの法則}です。\bkidx{あんぺーるのほうそく}{アンペールの法則}
\pheq{\oint_C\vec B\cdot\dd\vec r=\mu_0I_{\text{中}}}{T·m ＝ N/A$^2$ × A}
\label{ph:電-6:ampere}「閉じた道 $C$ にそって、磁場の道にそう部分×道の長さを1周たす（線積分 8-7）と、その道が囲む電流 $I_{\text{中}}$ の $\mu_0$ 倍になる」。道の形は自由に選べます。道を貫く電流の向きは、道を回る向きに右ねじを回したとき、ねじが進む向きを正とします。

\phboxhead{まっすぐな電流：円の道を選ぶ}
導線を中心とする半径 $r$ の円を道に選ぶと、対称性から、磁場は円にそっていて、円の上のどこでも同じ大きさ $B$ です。線積分は $B\times2\pi r$ なので、$B\times2\pi r=\mu_0I$。⑤と同じ $B=\mu_0I/(2\pi r)$ が、1行で出てきます。高校の式の $2\pi r$ は、円の道の長さだったのです。

\phboxhead{長いソレノイド}
\label{ph:電-6:solenoid}導線を、つつの周りにらせんの形にすき間なく巻いたコイルを\textbf{ソレノイド}といいます。\bkidx{それのいど}{ソレノイド}1 m あたり $n$ 回巻いた長いソレノイドに電流 $I$ を流すと、中では1つ1つの輪の磁場が同じ向きにそろって強め合い、つつの軸に平行な、一様な磁場ができます。外では、ほとんど打ち消し合って、磁場はとても弱くなります（長いほど0に近い）。

\begin{bkfigure}
\begin{tikzpicture}[x=1mm, y=1mm]
  \fill[bkBlue!8] (0,-10) rectangle (70,10);
  \foreach \xx in {2,6,...,66} {\sxgout{(\xx,11.5)}{1.5}\sxgin{(\xx,-11.5)}{1.5}}
  \foreach \yy in {-3,1.5,6} {\draw[sxg B] (8,\yy) -- (62,\yy);}
  \draw[bkRed!80!black, line width=1pt, dash pattern=on 3pt off 1.5pt] (24,-7.5) rectangle (44,22);
  \draw[bkRed!80!black, line width=1pt, -{Stealth[length=2mm]}] (30,-7.5) -- (38,-7.5);
  \node[sxg lab, anchor=east] at (23.4,-7.5) {A};
  \node[sxg lab, anchor=west] at (44.6,-7.5) {B};
  \node[sxg lab, anchor=west] at (44.6,22) {C};
  \node[sxg lab, anchor=east] at (23.4,22) {D};
  \node[sxg lab, anchor=south] at (34,22.4) {外：磁場ほぼ0};
  \node[sxg lab, anchor=north] at (34,-14.2) {道 ABCD（中の辺 AB の長さ $l$）};
  \node[anchor=west, font=\small, text width=50mm, align=left] at (76,6) {%
  $\odot$：巻いた導線の上側（電流は手前へ）。$\otimes$：下側（おくへ）。\par\vspace{2pt}
  道が囲む電流は、上の辺のはば $l$ の分の巻数 $nl$ 本ぶんで、$nlI$。};
\end{tikzpicture}
\bkfigcaption{長いソレノイドの断面と、アンペールの法則で選ぶ長方形の道。中の磁場は軸に平行で一様、外はほぼ0。}\label{ph:電-6:fig-solenoid}
\end{bkfigure}

\noindent 図\ref{ph:電-6:fig-solenoid}の長方形の道 ABCD を選びます。中の辺 AB（長さ $l$）では磁場は道にそって $B$ で、線積分は $Bl$。たての辺 BC・DA では磁場が道に垂直（外の部分は磁場0）なので0。外の辺 CD も磁場0なので0です。道が囲む電流は、巻いた導線 $nl$ 本ぶんの $nlI$ です。アンペールの法則から $Bl=\mu_0nlI$、だから
\pheq{B=\mu_0nI}{T ＝ N/A$^2$ × (1/m) × A}
です。高校の教科書の $H=nI$ に $\mu_0$ をかけたものです。中の磁場は、道の辺 AB をどの高さにとっても同じ答えになるので、\textbf{断面のどこでも同じ（一様）}です。つつの太さにもよりません。

\begin{phflow}[電流がつくる磁場]
\phflowrow{$\dd\vec B=\dfrac{\mu_0}{4\pi}\dfrac{I\,\dd\vec l\times\vec e_r}{r^2}$}{円の中心\\どの部分も距離 $a$}{$B=\dfrac{\mu_0I}{2a}$}
\phflowrow{$\oint_C\vec B\cdot\dd\vec r=\mu_0I_{\text{中}}$}{まっすぐな導線\\円の道を選ぶ}{$B=\dfrac{\mu_0I}{2\pi r}$}
\phflowrow{$\oint_C\vec B\cdot\dd\vec r=\mu_0I_{\text{中}}$}{長いソレノイド\\外の磁場0}{$B=\mu_0nI$}
\end{phflow}

\begin{phbaisuru}{1 m あたりの巻数 $n$}
中の磁場も2倍になります。電流を2倍にしても2倍です。電磁石を強くするには「たくさん巻いて、たくさん流す」。小学校の理科で調べた「電磁石の強さは、電流の大きさと巻数で変わる」は、この式のことです。磁場を決めるのは $nI$（1 m あたりの「巻数×電流」。単位 A/m）なので、工学では $NI$ を\textbf{アンペアターン}とよびます（電-7）。
\end{phbaisuru}

\begin{phmijika}[くぎの電磁石：小学校の実験を式で読む]
鉄のくぎにエナメル線を巻き、乾電池をつないで、くぎの先につくクリップの数を数える実験をしたことがあるでしょう。巻数を2倍にしたり、電池を2個直列にして電流をふやしたりすると、つくクリップがふえます。巻数を2倍にすると、同じ長さの所に2倍の導線が巻かれるので $n$ が2倍、電池を直列にすると電流 $I$ がふえます。どちらも $nI$ を大きくして、⑥の $B=\mu_0nI$ を強めているのです。鉄のくぎを入れると、空気のときより磁場がずっと強くなるわけは、電-7 で学びます。ただし、電池の数をふやしすぎると、エナメル線が熱くなります。実験は短い時間で行いましょう。
\end{phmijika}

\phboxhead{環状コイル（トロイダルコイル）}
ドーナツの形の芯に、導線を $N$ 回、一様に巻いたコイルを\textbf{環状コイル}といいます。中心から半径 $r$ の円を道に選ぶと、道はドーナツの中を通り、どの巻いた導線も1回ずつ貫くので、$B\times2\pi r=\mu_0NI$ です。
\pheq{B=\frac{\mu_0NI}{2\pi r}}{T}
ドーナツの外（中心の穴の中や、ずっと外）では、道が囲む電流の合計が0（行きと帰りが同じ数）なので、磁場は0です。磁場がコイルの中だけに閉じこめられるので、まわりの機械に影響しにくく、電源の部品によく使われます（電-7 の磁気回路の形）。

\begin{bkfigure}
\begin{tikzpicture}[x=1mm, y=1mm]
  \fill[bkBlue!8, even odd rule] (0,0) circle (20) (0,0) circle (11);
  \draw[bkInk, line width=0.6pt] (0,0) circle (20) (0,0) circle (11);
  \foreach \aa in {0,30,...,330} {\draw[bkOrange!80!black, line width=1.1pt] (\aa:9.5) -- (\aa:21.5);}
  \draw[bkRed!80!black, line width=1pt, dash pattern=on 3pt off 1.5pt] (0,0) circle (15.5);
  \draw[bkRed!80!black, line width=1pt, -{Stealth[length=2mm]}] (-60:15.5) arc (-60:-30:15.5);
  \node[sxg lab, anchor=west] at (24,0) {道（半径 $r$）};
  \draw[bkRed!80!black, line width=0.5pt] (22,-0.5) -- (15.6,-1.8);
  \node[sxg lab, anchor=north] at (0,-23) {環状コイル（$N$ 回巻き）};
  \begin{scope}[shift={(60,-8)}]
  \draw[bkBlue!80!black, line width=1.2pt, domain=0.6:1.4, samples=60] plot ({40*(\x-0.6)}, {16/\x});
  \draw[bkGray, line width=0.5pt, -{Stealth[length=2mm]}] (0,0) -- (38,0);
  \draw[bkGray, line width=0.5pt, -{Stealth[length=2mm]}] (0,0) -- (0,30);
  \node[sxg lab, anchor=north] at (36,-0.6) {$r$};
  \node[sxg lab, anchor=east] at (-0.6,27) {$B$};
  \node[sxg lab, anchor=north] at (2,-0.6) {内のはし};
  \node[sxg lab, anchor=south] at (32,13.8) {外のはし};
  \node[sxg lab, anchor=west] at (6,31) {内側ほど少し強い（$1/r$）};
  \end{scope}
\end{tikzpicture}
\bkfigcaption{環状コイルと、選ぶ円の道。中の磁場は $1/r$ に比例するので、ドーナツの内側ほど少し強い（例題4）。}\label{ph:電-6:fig-toroid}
\end{bkfigure}

\phboxhead{有限の長さのソレノイド：Python でたし合わせる}
本物のソレノイドは、長さに限りがあります。中の磁場は、どこでも $\mu_0nI$ なのでしょうか。半径\qty{2}{cm}、長さ\qty{10}{cm}、100回巻き（$n=1000$ /m）、電流\qty{1}{A}のソレノイドを、100個の円い電流の集まりとみなし、軸の上の磁場をたし合わせます。1つの円い電流（半径 $a$）が、軸の上で中心から $z$ 離れた所につくる磁場は、ビオ・サバールの法則を④と同じように積分して、$B=\mu_0Ia^2/\{2(a^2+z^2)^{3/2}\}$ です（$z=0$ で④の $\mu_0I/(2a)$）。
\begin{phcode}[Python]
import math
mu0 = 4 * math.pi * 1e-7
a, L, N, I = 0.02, 0.10, 100, 1.0     \# 半径・長さ・巻数・電流
zs = [-L/2 + L*(k + 0.5)/N for k in range(N)]   \# 100個の輪の位置
def B(z):                               \# 軸の上、中心から z の所
    return sum(mu0*I*a*a / (2*(a*a + (z - zk)**2)**1.5) for zk in zs)
for z in [0, 0.025, 0.05, 0.075, 0.10]:
    print(z, round(B(z)*1000, 3))       \# mT で表示
print('mu0 n I =', round(mu0*N/L*I*1000, 3))
\end{phcode}
\begin{phcodeout}
0 1.167
0.025 1.098
0.05 0.616
0.075 0.13
0.1 0.039
mu0 n I = 1.257
\end{phcodeout}
左から、中心からの距離〔m〕、磁場〔mT〕です。まん中では、無限に長いときの $\mu_0nI=\qty{1.257}{mT}$ の93\,\%、はし（$z=\qty{0.05}{m}$）ではその約半分です。はしで半分になるのは、はしから見ると、ソレノイドが片側にしかないからです。長さが半径の5倍あっても、はしの近くでは一様ではありません。

\begin{bkfigure}
\begin{tikzpicture}
\begin{axis}[bkplot, width=110mm, height=46mm, xmin=-0.1, xmax=0.1, ymin=0, ymax=1.45, xtick={-0.1,-0.05,0,0.05,0.1}, xticklabel style={/pgf/number format/fixed}, ytick={0.5,1.0}, yticklabels={0.5,1.0}, xlabel={軸の上の位置 $z$（m）}, ylabel={$B$（mT）}]
  \fill[bkOrange!12] (axis cs:-0.05,0) rectangle (axis cs:0.05,1.4);
  \addplot[bkcurve, domain=-0.1:0.1, samples=161] {1.2566370614*0.5*((0.05-x)/sqrt((0.05-x)^2+0.0004)+(0.05+x)/sqrt((0.05+x)^2+0.0004))};
  \addplot[bkGray, line width=1pt, dash pattern=on 4pt off 2pt, domain=-0.05:0.05] {1.2566};
  \fill[bkRed] (axis cs:0,1.167) circle (1.8pt) (axis cs:0.05,0.616) circle (1.8pt);
  \node[font=\sffamily\footnotesize, anchor=west, fill=white, inner sep=1pt] at (axis cs:0.053,1.26) {破線：無限に長いとき $\mu_0nI$};
  \node[font=\sffamily\footnotesize, anchor=west, text=bkRed!75!black] at (axis cs:0.054,0.7) {はしで約半分};
  \node[font=\sffamily\footnotesize, anchor=north] at (axis cs:-0.025,0.45) {コイルの中};
\end{axis}
\end{tikzpicture}
\bkfigcaption{長さ\qty{10}{cm}・半径\qty{2}{cm}のソレノイドの軸の上の磁場（色をつけた所がコイルの中）。点はプログラムの答え。}\label{ph:電-6:fig-finite}
\end{bkfigure}

\begin{phmodel}[長いソレノイド]
\phignore{巻いた導線のらせんのずれ（1回ごとに少し進む）。導線の太さ。}
\phassume{ソレノイドは半径に比べて十分長く、外の磁場は0。巻き方は一様。}
\phlimit{はしの近く（磁場は約半分になる）。長さが半径の数倍より短いとき。中に鉄を入れたとき（電-7）。}
\end{phmodel}

\bksubsection{7}{場所ごとの形：$\mathrm{rot}\,\vec B=\mu_0\vec j$ と $\mathrm{div}\,\vec B=0$}
\label{ph:電-6:sec7}
アンペールの法則は、閉じた道1周についての式（\textbf{積分形}）でした。道をとても小さくすると、場所ごとの式（\textbf{微分形}）になります。8-10⑪で確かめたように、まっすぐな導線の外では、磁場は導線の周りを回っているのに、回転（rot）は0でした。導線の中では、回転は電流密度 $\vec j$（断面 1 m$^2$ あたりの電流：電-4 ③）の $\mu_0$ 倍でした。ストークスの定理（8-11⑨）で道1周の線積分を面の積分に直すと、
\[\oint_C\vec B\cdot\dd\vec r=\iint_S\mathrm{rot}\,\vec B\cdot\vec n\,\dd S\qquad\text{と}\qquad\mu_0I_{\text{中}}=\iint_S\mu_0\vec j\cdot\vec n\,\dd S\]
が、どんな道・どんな面でも等しいので、
\pheq{\mathrm{rot}\,\vec B=\mu_0\vec j}{T/m ＝ N/A$^2$ × A/m$^2$}
\label{ph:電-6:rot}です。「\textbf{磁場のうずは、電流が流れている所にだけある}」と読みます。

もう1つの式は、①で見た「N極だけの磁石はない」「磁力線はいつも輪になる」ことです。どんな袋をとっても、入る磁力線と出る磁力線は同じ数なので、磁場の流束（8-8）は0です。8-9の発散のことばで書くと、
\pheq{\oiint_S\vec B\cdot\vec n\,\dd S=0\qquad\mathrm{div}\,\vec B=0}{T·m$^2$}
\label{ph:電-6:div}です。「\textbf{磁場には、わき出しも吸いこみもない}」と読みます。電場は電荷からわき出しました（$\mathrm{div}\,\vec E=\rho/\varepsilon_0$：電-2 ⑧）。磁場では、それにあたる「磁荷」が見つかっていません。

\begin{bkcenter}
\begin{bktable}{colspec={Q[l,m,wd=26mm] Q[l,m,wd=44mm] Q[l,m,wd=44mm]}, rowsep=1.5pt, row{1}={bg=bkTheme!10, font=\sffamily\small}, cells={font=\small}}
 & 止まった電荷の電場（電-2） & 止まった電流の磁場（この章） \\
わき出し（div） & $\mathrm{div}\,\vec E=\rho/\varepsilon_0$（電荷から出る） & $\mathrm{div}\,\vec B=0$（わき出しなし） \\
うず（rot） & $\mathrm{rot}\,\vec E=\vec0$（うずなし） & $\mathrm{rot}\,\vec B=\mu_0\vec j$（電流の所にうず） \\
力線の形 & ＋から出て－に入る & いつも閉じた輪 \\
\end{bktable}
\bktabcaption{電場と磁場の「わき出し」と「うず」。時間で変わらないときの形。時間で変わると、電場にもうずができる（電-8）。}\label{ph:電-6:tab-divrot}
\end{bkcenter}

\begin{seikaku}
アンペールの法則 $\oint\vec B\cdot\dd\vec r=\mu_0I_{\text{中}}$ は、電流が時間で変わらないときの形です。コンデンサーに電流を流しこんでいる途中のように、電流がとちゅうで途切れているように見える所では、このままでは矛盾が起きます。マクスウェルは、電場が時間で変わる所にも磁場のうずができる項を足して、この矛盾をなくしました（8-14④、電-14）。
\end{seikaku}

\begin{motto}[4]
\phboxhead{ベクトルポテンシャル}
$\mathrm{div}\,\vec B=0$ の場は、別のベクトルの場 $\vec A$ の回転として $\vec B=\mathrm{rot}\,\vec A$ と書けます（8-10⑩の $\mathrm{div}\,(\mathrm{rot}\,\vec F)=0$ の逆の読み方）。この $\vec A$ を\textbf{ベクトルポテンシャル}といいます。電場が電位 $\phi$ の勾配で書けたように（電-2 ⑦）、磁場はベクトルポテンシャルの回転で書けます。電流の小さな部分は、同じ向きの $\vec A$ を $\mu_0I\,\dd\vec l/(4\pi r)$ だけつくり（距離に反比例。電位 $kq/r$ と同じ形）、その回転をとると④のビオ・サバールの法則が出てきます。電波を出すアンテナの計算（電-15）や、量子力学では、$\vec A$ が主役になります。詳しくは電磁気学の専門の本で学びます。
\end{motto}

\bksubsection{8}{電流が磁場から受ける力}
\label{ph:電-6:sec8}
導線の中では、たくさんの電荷が動いています。磁場の中の導線は、その電荷1つ1つが受けるローレンツ力を合わせた力を受けます。長さ $l$、断面積 $S$ の導線の中に、電荷 $q$ の粒が1 m$^3$ あたり $n_{\mathrm c}$ 個あり、速さ $v$ で流れているとします。導線の中の粒の数は $n_{\mathrm c}Sl$ 個、1個が受ける力は $qvB\sin\theta$ なので、合わせると $n_{\mathrm c}Sl\times qvB\sin\theta$ です。電-4 ②で学んだように、電流は $I=n_{\mathrm c}qvS$ なので、
\pheq{F=IBl\sin\theta\qquad\vec F=I\,\vec l\times\vec B}{N ＝ A × T × m}
\label{ph:電-6:force}です（$\vec l$ は電流の向きの、長さ $l$ のベクトル。4-14⑧の $\vec L$ と同じ）。$\theta$ は電流と磁場の間の角です。\bkidx{でんりゅうがうけるちから}{電流が磁場から受ける力}

\noindent\begin{minipage}[c]{0.46\linewidth}\centering
\begin{tikzpicture}[x=1mm, y=1mm]
  \foreach \xx in {4,14,...,54} {\foreach \yy in {-2,18,28} {\sxgin{(\xx,\yy)}{1.6}}}
  \draw[bkOrange!80!black, line width=2pt] (-2,8) -- (60,8);
  \draw[sxg I] (22,4.5) -- (34,4.5);
  \node[sxg lab, anchor=west] at (35,4.5) {電流 $I$};
  \draw[sxg F] (30,9.5) -- (30,26);
  \node[sxg lab, anchor=west] at (31.5,23) {力 $\vec F$};
  \node[sxg lab, anchor=north] at (30,-5.5) {磁場 $\vec B$ はおくへ（$\otimes$）};
\end{tikzpicture}
\end{minipage}\hfill
\begin{minipage}[c]{0.51\linewidth}\raggedright
磁場が紙のおくへ（$\otimes$）、電流が右へなら、$\vec l\times\vec B$ は上向きです（右向き×おく向き＝上向き）。中学で学んだフレミングの左手の法則（中指＝電流、人さし指＝磁場、親指＝力）でも、同じ向きになります。\par\vspace{2pt}
4-14⑧の\qty{3}{A}・\qty{0.2}{m}・\qty{0.5}{T}なら、$3\times0.5\times0.2=\qty{0.30}{N}$ です（例題5）。
\end{minipage}

\par\vspace{4pt}
\begin{phyomi}[$F=IBl\sin\theta$]
「電流と磁場が直角のとき、力がいちばん大きく、平行なら0」。力の向きは、電流にも磁場にも垂直です。電流を反対にすると、力も反対になります。モーター（電-12）は、この力でコイルを回します。スピーカーは、この力で、磁石の中のコイルとそれにつけた紙の振動板を前後にゆすります。
\end{phyomi}

\begin{phmijika}[スピーカーとイヤホン]
スピーカーの奥には、円い溝のある磁石があり、溝の中では磁場が中心から外へ向いています（放射状）。溝にコイルが入っていて、コイルの導線はどこでも磁場と直角です。音の信号の電流が流れると、コイル全体が前か後ろへ押され、電流の向きが変わるたびに押す向きが変わって、振動板が空気をふるわせます。
\end{phmijika}

\bksubsection{9}{平行な電流の間の力と、アンペアの定義}
\label{ph:電-6:sec9}
平行な2本の長い導線に、電流 $I_1$・$I_2$ を流します。間の距離を $r$ とすると、導線1が導線2の所につくる磁場は $B_1=\mu_0I_1/(2\pi r)$ で、導線2に直角です。導線2の長さ $l$ の部分は、⑧の力 $I_2B_1l$ を受けます。
\pheq{F=\frac{\mu_0I_1I_2\,l}{2\pi r}}{N ＝ N/A$^2$ × A × A × m÷m}
\label{ph:電-6:parallel}向きを右ねじと⑧の外積で調べると、\textbf{同じ向きの電流は引き合い、反対向きの電流は退け合い}ます。導線1も、導線2の磁場から、同じ大きさで反対向きの力を受けます（力-2 の作用・反作用）。\bkidx{へいこうでんりゅう}{平行電流の間の力}

\begin{bkfigure}
\begin{tikzpicture}[x=1mm, y=1mm]
  \draw[bkOrange!80!black, line width=2pt] (0,0) -- (0,34);
  \draw[bkOrange!80!black, line width=2pt] (24,0) -- (24,34);
  \draw[sxg I] (-4,3) -- (-4,15);
  \draw[sxg I] (28,3) -- (28,15);
  \node[sxg lab, anchor=east] at (-5,9) {$I_1$};
  \node[sxg lab, anchor=west] at (29,9) {$I_2$};
  \sxgin{(12,10)}{1.6}\sxgin{(18,10)}{1.6}
  \node[sxg lab, anchor=north] at (15,7.5) {$I_1$ の磁場};
  \draw[sxg F] (23,24) -- (14,24);
  \draw[sxg F] (1,24) -- (10,24);
  \node[sxg lab, anchor=north] at (12,-1) {同じ向き：引き合う};
  \begin{scope}[shift={(70,0)}]
  \draw[bkOrange!80!black, line width=2pt] (0,0) -- (0,34);
  \draw[bkOrange!80!black, line width=2pt] (24,0) -- (24,34);
  \draw[sxg I] (-4,3) -- (-4,15);
  \draw[sxg I] (28,15) -- (28,3);
  \node[sxg lab, anchor=east] at (-5,9) {$I_1$};
  \node[sxg lab, anchor=west] at (29,9) {$I_2$};
  \draw[sxg F] (25,24) -- (36,24);
  \draw[sxg F] (-1,24) -- (-12,24);
  \node[sxg lab, anchor=north] at (12,-1) {反対向き：退け合う};
  \end{scope}
\end{tikzpicture}
\bkfigcaption{平行な電流の間の力。電流の矢印は細い実線、力の矢印は二重線。}\label{ph:電-6:fig-parallel}
\end{bkfigure}

\phboxhead{アンペアの定義の移り変わり}
1820年の秋、フランスのアンペールは、エルステッドの発見を聞いてすぐに、平行な電流が引き合ったり退け合ったりすることを実験で示しました。この力は、電流の大きさを決める物差しになります。1948年から2019年5月19日まで、電流の単位アンペアは、次のように決められていました。
\begin{bkcenter}
\begin{minipage}{0.88\linewidth}\small
真空の中に、\qty{1}{m}離して平行に置いた、無限に長く、とても細い2本の導線に同じ電流を流したとき、導線 1 m あたり $2\times10^{-7}$\,N の力がはたらく電流を\qty{1}{A}とする。
\end{minipage}
\end{bkcenter}
上の式に $I_1=I_2=\qty{1}{A}$、$r=\qty{1}{m}$、$l=\qty{1}{m}$ を入れると、$\mu_0/(2\pi)=2\times10^{-7}$、つまり $\mu_0=4\pi\times10^{-7}$ になります（例題6）。$\mu_0$ がちょうどこの値だったのは、アンペアをこう決めたからです。けれども、「無限に長い導線」は作れず、力を精密にはかるのもたいへんです。2019年5月20日からは、電気素量 $e$ の値を決まった数にして、「\qty{1}{A}は、1秒に $1/(1.602\,176\,634\times10^{-19})$ 個の電気素量が流れる電流」と決めています（③）。

\bksubsection{10}{磁場の中の荷電粒子：円運動・速度選別器・ホール効果}
\label{ph:電-6:sec10}
\phboxhead{磁場に直角に入った粒子は、円を回る}
\label{ph:電-6:lorentz}一様な磁場 $\vec B$ の中に、質量 $m$、電荷 $q$ の粒子が、磁場に直角な速度 $\vec v$ で入ります。ローレンツ力の大きさは $f=qvB$ で、向きはいつも速度に垂直です。力が速度に垂直なので、速さ $v$ は変わらず、向きだけが一定の割合で曲がります。これは\textbf{等速円運動}です（力-1 ⑤）。中心へ向かう力 $qvB$ が、円運動の運動方程式（力-2 ⑦）の向心力になるので、
\[m\frac{v^2}{r}=qvB\]
\pheq{r=\frac{mv}{qB}}{m ＝ kg × m/s ÷ (C × T)}
です。1周する時間（周期）は $2\pi r/v=2\pi m/(qB)$ で、1秒に回る回数（周波数）は
\pheq{f=\frac{qB}{2\pi m}}{Hz ＝ C × T ÷ kg}
です。\textbf{速さによらない}のが大事な点です。速い粒子は大きな円を、遅い粒子は小さな円を回り、1周の時間は同じです。\bkidx{さいくろとろんしゅうはすう}{サイクロトロン周波数}

\begin{bkfigure}
\begin{tikzpicture}[x=1mm, y=1mm]
  \foreach \xx in {-24,-12,...,24} {\foreach \yy in {-24,-12,12,24} {\sxgout{(\xx,\yy)}{1.3}}}
  \draw[bkGreen!50!black, line width=1.2pt, dash pattern=on 4pt off 2pt] (0,0) circle (18);
  \fill[bkRed] (18,0) circle (1.6);
  \node[sxg lab, anchor=south west, fill=white, inner sep=1pt] at (19,0.8) {$+q$};
  \draw[sxg I] (18,0) -- (18,-14);
  \node[sxg lab, anchor=east, fill=white, inner sep=1pt] at (16.6,-8.5) {$\vec v$};
  \draw[sxg F] (18,0) -- (6,0);
  \node[sxg lab, anchor=south, fill=white, inner sep=1pt] at (10,0.8) {$q\vec v\times\vec B$};
  \node[sxg lab, anchor=north] at (0,-28) {磁場は手前へ（$\odot$）};
  \node[anchor=west, font=\small, text width=52mm, align=left] at (32,4) {%
  正の電荷は、磁場が手前へ向かうとき、手前から見て\textbf{右回り}（時計回り）に回る（負の電荷は左回り）。\par\vspace{2pt}
  半径 $r=\dfrac{mv}{qB}$：速いほど・重いほど大きく、磁場が強いほど小さい。\par\vspace{2pt}
  1周の時間 $\dfrac{2\pi m}{qB}$ は、速さによらない。};
\end{tikzpicture}
\bkfigcaption{一様な磁場の中の荷電粒子の円運動（破線が通り道）。}\label{ph:電-6:fig-cyclo}
\end{bkfigure}

\begin{phflow}[ローレンツ力から]
\phflowrow{$\vec F=q\vec v\times\vec B$}{導線の中の電荷を\\全部たす（$I=n_{\mathrm c}qvS$）}{$F=IBl\sin\theta$}
\phflowrow{$F=IBl$}{となりの導線の\\磁場 $\mu_0I_1/(2\pi r)$}{$F=\dfrac{\mu_0I_1I_2l}{2\pi r}$}
\phflowrow{$\vec F=q\vec v\times\vec B$}{$\vec v\perp\vec B$\\一様な磁場}{$r=\dfrac{mv}{qB}$}
\end{phflow}

\begin{phbaisuru}{速さ $v$}
半径は2倍になります。周期と周波数は変わりません。磁場を2倍にすると、半径は半分、周波数は2倍です。
\end{phbaisuru}

\phboxhead{サイクロトロンと、比電荷の測定}
円を回る周波数が速さによらないことを使った加速器が\textbf{サイクロトロン}です。D字形の2つの金属の箱を向かい合わせに置き、その間のすき間に、粒子の回る周波数と同じ周波数の電圧をかけます。粒子は、すき間を通るたびに加速され、少しずつ大きな円を描きながら速くなります。陽子なら、\qty{1}{T}の磁場で約\qty{15}{MHz}です（例題9）。また、加速の電圧と円の半径をはかると、⑩の式から電荷と質量の比 $q/m$（比電荷）が求められます。電子の比電荷は、物理シリーズの『原子』（原-1）で、この考えではかります。

\phboxhead{速度選別器：電場と磁場で、決まった速さの粒子だけを通す}
電場 $\vec E$ と磁場 $\vec B$ を、たがいに直角にかけた所に、両方に直角な向きに粒子を入れます。正の電荷は、電場から $qE$、磁場から $qvB$ の力を、反対向きに受けるように向きをそろえておきます。2つがつり合う $qE=qvB$ の粒子、つまり速さ
\pheq{v=\frac{E}{B}}{m/s ＝ (V/m) ÷ T}
の粒子だけが、まっすぐ進みます。これより速い粒子は磁場の力が勝って、遅い粒子は電場の力が勝って、曲がってしまいます（例題13）。電荷や質量によらず、速さだけで選べるのが特ちょうです。

\phboxhead{ホール効果：磁場をはかる素子}
板の形の導体や半導体に、長さの向きに電流 $I$ を流し、板に垂直に磁場 $B$ をかけます。中を流れる電荷は、ローレンツ力で板の横の面へ押されて、たまります。たまった電荷は、板の横の向きに電場 $E$ をつくり、その電場の力 $qE$ が、ローレンツ力 $qvB$ とつり合うと、それ以上たまりません。このとき、板の横のはしとはしの間に電圧 $V_{\mathrm H}$ が出ます。これを\textbf{ホール効果}といいます（1879年、アメリカのホール）。\bkidx{ほーるこうか}{ホール効果}

\noindent\begin{minipage}[c]{0.44\linewidth}\centering
\begin{tikzpicture}[x=1mm, y=1mm]
  \filldraw[fill=bkGray!15, draw=bkInk, line width=0.7pt] (0,0) -- (40,0) -- (52,10) -- (12,10) -- cycle;
  \filldraw[fill=bkGray!30, draw=bkInk, line width=0.7pt] (0,0) -- (40,0) -- (40,-3) -- (0,-3) -- cycle;
  \draw[sxg I] (-9,-1.5) -- (-1,-1.5);
  \node[sxg lab, anchor=south] at (-5,-1) {$I$};
  \draw[sxg B, line width=1.3pt] (6,-13) -- (6,-4.5);
  \draw[sxg B, line width=1.3pt] (14,12) -- (14,22);
  \node[sxg lab, anchor=west] at (15,19) {$B$（上へ）};
  \node[sxg lab, anchor=north] at (24,-4.5) {＋＋＋（手前の面）};
  \node[sxg lab, anchor=south] at (36,10.6) {－－－（おくの面）};
\end{tikzpicture}
\end{minipage}\hfill
\begin{minipage}[c]{0.53\linewidth}\raggedright
正の電荷が流れる板（はば $w$、厚さ $t_{\mathrm H}$）では、$qE=qvB$、$E=V_{\mathrm H}/w$、電流 $I=n_{\mathrm c}qv\times wt_{\mathrm H}$ から、$v$ と $w$ を消して、
\[V_{\mathrm H}=\frac{IB}{n_{\mathrm c}qt_{\mathrm H}}\]
です（$n_{\mathrm c}$ は流れる電荷の数密度）。電圧は $B$ に比例するので、$V_{\mathrm H}$ から磁場がわかります。$n_{\mathrm c}$ が小さいほど電圧が大きいので、ホール素子には、金属より電荷の少ない半導体を使います。
\end{minipage}

\par\vspace{4pt}
\noindent 図の向きでは、正の電荷は手前の面に集まります。流れているのが負の電荷（電子）なら、電子は電流と反対向きに動くので、ローレンツ力は同じ手前向きになり、手前の面が\textbf{負}になります。ホール電圧の符号から、電流を運んでいるのが正の電荷か負の電荷かがわかるのです。スマホの電子コンパス、モーターの回転の位置を読むセンサー、導線を切らずに電流をはかる電流センサー（電流の周りの磁場をはかる：⑤）に使われています。

\bksubsection{11}{磁場の中のコイルにはたらくトルク}
\label{ph:電-6:sec11}
一様な磁場 $B$ の中に、長方形のコイル（辺の長さ $a$・$b$、面積 $S=ab$、巻数 $N$）を置き、電流 $I$ を流します。コイルは、辺 $b$ を通る軸の周りに回れるとします。磁場に直角な辺 $a$ には、⑧の力 $NIBa$ が、2本の辺で反対向きにはたらきます。2つの力は同じ線の上にないので、コイルを回そうとします（力のモーメント：4-14・力-6）。

\begin{bkfigure}
\begin{tikzpicture}[x=1mm, y=1mm]
  \foreach \yy in {-6,6,18} {\draw[sxg B] (-30,\yy) -- (30,\yy);}
  \node[sxg lab, anchor=west] at (31,18) {$\vec B$};
  \draw[bkOrange!80!black, line width=2pt] (-9.64,-11.49) -- (9.64,11.49);
  \sxgout{(9.64,11.49)}{2}
  \sxgin{(-9.64,-11.49)}{2}
  \fill[bkInk] (0,0) circle (0.9);
  \node[sxg lab, anchor=south east, fill=white, inner sep=1pt] at (-1.5,1.2) {回る軸};
  \draw[sxg F] (9.64,11.49) ++(0,2.5) -- ++(0,9);
  \draw[sxg F] (-9.64,-11.49) ++(0,-2.5) -- ++(0,-9);
  \node[sxg lab, anchor=west, fill=white, inner sep=1pt] at (11.5,22.5) {$NIBa$};
  \node[sxg lab, anchor=east] at (-11.5,-20) {$NIBa$};
  \draw[bkGreen!50!black, line width=1pt, -{Stealth[length=2mm]}] (-60:6) arc (-60:-10:6);
  \draw[bkInk, line width=0.5pt, dash pattern=on 1pt off 1pt] (0,0) -- (12.3,-10.3);
  \node[sxg lab, anchor=north] at (13,-11) {面に垂直な向き};
  \node[anchor=west, font=\small, text width=56mm, align=left] at (44,0) {%
  上から（辺 $b$ の向きから）見た図。$\odot$ と $\otimes$ は、2本の辺 $a$ を流れる電流。コイルの面に垂直な向きと磁場の間の角を $\theta$ とすると、うでの長さは $\tfrac b2\sin\theta$。\par\vspace{2pt}
  2つの力のモーメントの和：\par
  $T=2\times NIBa\times\tfrac b2\sin\theta$\par
  $\phantom{T}=NISB\sin\theta$};
\end{tikzpicture}
\bkfigcaption{磁場の中の長方形のコイル。2本の辺が反対向きに押され、コイルを回す。}\label{ph:電-6:fig-torque}
\end{bkfigure}

\noindent コイルの面に垂直な向き（電流の向きに右ねじを回したときに進む向き）と、磁場の間の角を $\theta$ とすると、トルク（回そうとする働き）の大きさは
\pheq{T=NISB\sin\theta}{N·m ＝ A × m$^2$ × T}
\label{ph:電-6:torque}です。\bkidx{こいるのとるく}{コイルにはたらくトルク}コイルの面が磁場と平行（$\theta=90^\circ$）のとき、いちばん大きく、面が磁場と直角（$\theta=0$）になると0になります。コイルは、面に垂直な向きが磁場の向きにそろうように回ります。小さな円い電流が平たい磁石のようにふるまう（④）のと同じで、$NIS$ が「コイルという磁石の強さ」を表しています（方位磁針の針がそろうのと同じ）。

\begin{phmijika}[針の電流計とモーター]
アナログの電流計では、ばねにつないだ小さなコイルが磁石の間にあり、電流に比例するトルクで針が回ります。ばねが戻す力と、トルクがつり合う角が、電流の大きさを指します。$\theta=0$ で回るのをやめてしまうと、モーターになりません。そこで直流モーターでは、コイルが半回転するごとに電流の向きを切りかえる部品（整流子）をつけて、同じ向きに回し続けます（電-12 ③）。
\end{phmijika}

\bksubsection{12}{地球の磁場：偏角と伏角}
\label{ph:電-6:sec12}
地球の磁場は、北の向きに水平なわけではありません。日本では、磁場は北へ向かいながら、地面にもぐりこむように斜め下を向いています。水平な面と磁場の間の角を\textbf{伏角}、真北（地図の北）と、磁場の水平の部分の向き（方位磁針が指す北）の間の角を\textbf{偏角}といいます。\bkidx{ふっかく}{伏角}\bkidx{へんかく}{偏角}

\begin{bkfigure}
\begin{tikzpicture}[x=1mm, y=1mm]
  \fill[bkGreen!15] (-4,-2) rectangle (44,0);
  \draw[bkInk, line width=0.6pt] (-4,0) -- (44,0);
  \draw[bkBlue!80!black, line width=1.6pt, -{Stealth[length=3mm]}] (0,22) -- (30,-13);
  \draw[bkBlue!60!black, line width=1pt, dash pattern=on 3pt off 1.5pt, -{Stealth[length=2mm]}] (0,22) -- (30,22);
  \draw[bkBlue!60!black, line width=1pt, dash pattern=on 3pt off 1.5pt, -{Stealth[length=2mm]}] (30,22) -- (30,-13);
  \node[sxg lab, anchor=south] at (15,22.4) {水平 約 $3.0\times10^{-5}$\,T（北へ）};
  \node[sxg lab, anchor=west] at (31,5) {下向き 約 $3.5\times10^{-5}$\,T};
  \node[sxg lab, anchor=east] at (9,4) {全体 約 $4.6\times10^{-5}$\,T};
  \draw[bkInk, line width=0.5pt] (8,22) arc (0:-49:8);
  \node[sxg lab, anchor=west] at (8.4,17.5) {伏角 約49°};
  \node[sxg lab, anchor=north] at (20,-13) {横から見る（東京）};
  \begin{scope}[shift={(84,6)}]
  \draw[bkInk, line width=0.8pt, -{Stealth[length=2.4mm]}] (0,0) -- (0,22);
  \node[sxg lab, anchor=south] at (0,22.4) {真北（地図）};
  \draw[bkRed!80!black, line width=1.2pt, -{Stealth[length=2.4mm]}] (0,0) -- (-2.8,21.8);
  \node[sxg lab, anchor=east] at (-3.6,15) {方位磁針};
  \draw[bkInk, line width=0.5pt] (0,12) arc (90:97.3:12);
  \node[sxg lab, anchor=west] at (1.5,12) {偏角 約7°（西）};
  \node[sxg lab, anchor=north] at (0,-2) {上から見る};
  \end{scope}
\end{tikzpicture}
\bkfigcaption{日本（東京あたり）の地球の磁場。数は丸めた目安で、場所と年で少しずつ変わる。}\label{ph:電-6:fig-earth}
\end{bkfigure}

\noindent 東京あたりでは、磁場全体の強さは約 $4.6\times10^{-5}$\,T、伏角は約49°、水平の部分は約 $3.0\times10^{-5}$\,T です。偏角は西に約7°で、方位磁針の北は、真北より少し西を指します（北海道では約9°、沖縄では約5°と、北ほど大きい）。地図とコンパスで山を歩くときは、この偏角を直してから方角を読みます。伏角は北へ行くほど大きく、地球の磁石のS極の真上（北極の近くの北極海の上。少しずつ動いている）では90°、赤道の近くでは0°に近くなります。方位磁針の針は、ふつう、水平に回る部分だけを見ています。

\begin{phmada}[渡り鳥は、地球の磁場をどうやって感じている？]
渡り鳥の多くは、地球の磁場の向きを感じて、何千 km も先の方角を知っています。鳥を磁場を変えた部屋に入れると、飛び立とうとする向きが変わることが、実験で確かめられています。そのしくみとして有力なのが、目の中のたんぱく質（クリプトクロム）に光が当たってできる、ペアになった電子の状態が、地球の磁場の向きで少し変わる、という考えです。2021年には、ヨーロッパの渡り鳥（ヨーロッパコマドリ）のクリプトクロム4というたんぱく質が、試験管の中で磁場に反応することが報告されました。けれども、その実験の磁場は地球の磁場より強く、生きた鳥の目の中で本当にそれが起き、脳に方角を伝えているのかは、まだ確かめられていません。地球の磁場は $5\times10^{-5}$\,T くらいととても弱く、分子の中のエネルギーを少し変えるだけです。それを、どのように確かな感覚にしているのかは、生き物の物理の大きななぞの1つです。
\end{phmada}

\begin{phsaki}{磁場の閉じこめと加速器}
この章では、一様な磁場の中の円運動を調べました。磁場の強さが場所で変わると、粒子は円を回りながら、ゆっくり横へずれていきます（ドリフト）。核融合の研究では、1億度をこえるプラズマ（電子と原子核がばらばらになった気体）を、ドーナツの形の強い磁場の中に閉じこめます。粒子の加速器では、光の速さに近い粒子の運動を、相対論（原-11）を入れて計算します。どちらも、プラズマ物理や加速器の専門の本で学びます。
\end{phsaki}



\begin{taisetsu}
\def\tsl#1{\makebox[7.6\zw][l]{\sffamily\bfseries\small\color{\bkTT}#1}}
\tsl{ローレンツ力}$\vec F=q\vec v\times\vec B$（磁場を決める式。単位 T＝N/(A·m)）\\[3pt]
\tsl{ビオ・サバール}$\dd\vec B=\dfrac{\mu_0}{4\pi}\dfrac{I\,\dd\vec l\times\vec e_r}{r^2}$、$\mu_0=4\pi\times10^{-7}$ N/A$^2$\\[3pt]
\tsl{アンペール}$\oint_C\vec B\cdot\dd\vec r=\mu_0I_{\text{中}}$、$\mathrm{rot}\,\vec B=\mu_0\vec j$、$\mathrm{div}\,\vec B=0$\\[3pt]
\tsl{磁場の式}直線 $\dfrac{\mu_0I}{2\pi r}$、円の中心 $\dfrac{\mu_0I}{2a}$、ソレノイド $\mu_0nI$、環状 $\dfrac{\mu_0NI}{2\pi r}$\\[3pt]
\tsl{力}導線 $F=IBl\sin\theta$、平行電流 $\dfrac{\mu_0I_1I_2l}{2\pi r}$、コイル $T=NISB\sin\theta$\\[3pt]
\tsl{粒子}$r=\dfrac{mv}{qB}$、$f=\dfrac{qB}{2\pi m}$、速度選別 $v=\dfrac EB$、ホール $V_{\mathrm H}=\dfrac{IB}{n_{\mathrm c}qt_{\mathrm H}}$
\end{taisetsu}

\begin{kotoba}
\small
\bkword{磁場（磁束密度）}{動く電荷に力をおよぼす空間の性質。ローレンツ力 $q\vec v\times\vec B$ で決める。単位 T。}
\bkword{右ねじの法則}{電流の向きに右ねじを進めるとき、ねじを回す向きが磁場の向き。}
\bkword{真空の透磁率}{電流がつくる磁場の式に入る定数 $\mu_0=4\pi\times10^{-7}$ N/A$^2$（2019年からは測る値）。}
\bkword{ビオ・サバールの法則}{電流の小さな部分がつくる磁場の式。たし合わせて全体の磁場を求める。}
\bkword{アンペールの法則}{閉じた道にそった磁場の線積分は、道が囲む電流の $\mu_0$ 倍。}
\bkword{ホール効果}{電流と磁場に垂直な向きに電圧が出ること。磁場をはかるのに使う。}
\end{kotoba}

\bkkigocards{\bkkigocard{$\vec B$}{ビー}{磁場（磁束密度）〔T〕。高校の教科書の「磁束密度」。真空では高校の「磁場の強さ」$H$ の $\mu_0$ 倍（$H$ は電-7 から）}{magnetic flux density}{電-6}}
{\bkkigocard{$\mu_0$}{ミューゼロ}{真空の透磁率 $4\pi\times10^{-7}$ N/A$^2$。物質の透磁率 $\mu$ は電-7。第1巻の摩擦係数 $\mu$・第3巻の線密度 $\mu$ とは別の量}{magnetic constant}{電-6}}

%% ======================== 例題 ========================
\Needspace{10\baselineskip}\bksection{reidai}{例題でたしかめよう}
\label{ph:電-6:ex1}
\begin{phreidai}[方位磁針はどれだけふれる？（はてな？）]
南北に張った長い導線の真下\qty{5.0}{cm}の所に、方位磁針を置きました。導線に南から北へ\qty{10}{A}の電流を流すと、針は北から何度ふれますか。地球の磁場の水平の部分は北向きに $3.0\times10^{-5}$\,T とします。
\phbamen
方位磁針は、地球の磁場と電流の磁場を合わせた磁場の向きを向きます。2つの磁場のベクトルの和を求める場面です。
\phmushi
電池やつないだ導線のほかの部分の電流の磁場は無視します。導線は十分長い、まっすぐな導線とします。針は水平にしか回らないので、磁場のたての部分は考えません。
\phzu
図\ref{ph:電-6:fig-hatena}。地球の磁場 $B_{\text{地}}=3.0\times10^{-5}$\,T（北向き）、電流の磁場 $B$（東西の向き）、ふれの角 $\theta$。
\phriyu
長いまっすぐな電流の磁場は $B=\mu_0I/(2\pi r)$（⑤）。向きは右ねじの法則で、導線の真下では東西の向き。北向きと東西向きは直角なので、$\tan\theta=B/B_{\text{地}}$。
\phtateru
\[B=\frac{\mu_0I}{2\pi r},\qquad\tan\theta=\frac{B}{B_{\text{地}}}\]
\phtoku
\[B=\frac{4\pi\times10^{-7}\times10}{2\pi\times0.050}=\frac{2\times10^{-7}\times10}{0.050}=4.0\times10^{-5}\ \text{T}\]
\[\tan\theta=\frac{4.0\times10^{-5}}{3.0\times10^{-5}}\approx1.33,\qquad\theta\approx53^\circ\]
向きは、電流を手前に向けて見ると導線の真下で磁場が西向き（②）なので、西へふれます。
\phtashikame
単位：$\dfrac{\text{N/A}^2\times\text{A}}{\text{m}}=\dfrac{\text{N}}{\text{A·m}}=\text{T}$。\\
大きさ：8-10⑪の数と同じ $4\times10^{-5}$\,T で、地球の磁場と同じくらいなので、大きくふれるのは自然です。\\
極端な場合：電流が0なら $\theta=0$（北を向く）。電流をいくら大きくしても、$\tan\theta$ が大きくなるだけで、90°にはなりません。
\kotae 西へ約53°
\end{phreidai}
\begin{jibun}
例題1で、（1）電流を\qty{5.0}{A}にすると、何度ふれますか。（2）\qty{10}{A}のままで、ちょうど45°ふれるのは、導線から何 cm の所ですか。
\bkjans{（1）$B=2.0\times10^{-5}$\,T、$\tan\theta=0.67$ で約34°。（2）$B=3.0\times10^{-5}$\,T になる所。$r=2\times10^{-7}\times10\div(3.0\times10^{-5})$ で約\qty{6.7}{cm}}
\end{jibun}

\label{ph:電-6:ex2}
\begin{phreidai}[円いコイルの中心の磁場]
半径\qty{5.0}{cm}の円いコイルを10回巻き、\qty{1.0}{A}の電流を流します。中心の磁場の強さを求めましょう。
\phbamen
円い電流の中心の磁場を求める場面です。
\phmushi
10回の輪は、どれも同じ位置に重なっているとみなします（巻いた幅と導線の太さは無視）。引き出し線の電流の磁場も無視します。
\phzu
図\ref{ph:電-6:fig-circle}。$a=\qty{0.050}{m}$、$N=10$、$I=\qty{1.0}{A}$。
\phriyu
1つの輪の中心の磁場は $\mu_0I/(2a)$（④）。10回巻くと、同じ向きの磁場が10個重なるので10倍。
\phtateru
\[B=\frac{\mu_0NI}{2a}\]
\phtoku
\[B=\frac{4\pi\times10^{-7}\times10\times1.0}{2\times0.050}=4\pi\times10^{-5}\approx1.3\times10^{-4}\ \text{T}\]
\phtashikame
単位：N/A$^2$×A÷m＝T。\\
大きさ：地球の磁場（約 $4.6\times10^{-5}$\,T）の約3倍です。方位磁針をコイルの中心に置けば、はっきりふれます。\\
極端な場合：$a$ を大きくすると0に近づきます（電流が遠くなる）。
\kotae 約 $1.3\times10^{-4}$\,T
\end{phreidai}
\begin{jibun}
半径\qty{10}{cm}、50回巻きの円いコイルに\qty{0.50}{A}の電流を流します。中心の磁場を求めましょう。
\bkjans{$4\pi\times10^{-7}\times50\times0.50\div(2\times0.10)$ で、約 $1.6\times10^{-4}$\,T}
\end{jibun}

\label{ph:電-6:ex3}
\begin{phreidai}[ソレノイドの中の磁場]
1 m あたり1000回巻いた長いソレノイドに、\qty{2.0}{A}の電流を流します。中の磁場の強さを求めましょう。
\phbamen
長いソレノイドの中の一様な磁場を求める場面です。
\phmushi
ソレノイドは十分長いとし、はしの近くは考えません。中は空気（真空と同じ）とします。
\phzu
図\ref{ph:電-6:fig-solenoid}。$n=1000$ /m、$I=\qty{2.0}{A}$。
\phriyu
アンペールの法則で長方形の道を選ぶと、$B=\mu_0nI$（⑥）。
\phtateru
\[B=\mu_0nI\]
\phtoku
\[B=4\pi\times10^{-7}\times1000\times2.0\approx2.5\times10^{-3}\ \text{T}\]
\phtashikame
単位：N/A$^2$×(1/m)×A＝N/(A·m)＝T。\\
大きさ：地球の磁場の約50倍。中に鉄の芯を入れると、さらに数百〜数千倍になります（電-7）。\\
極端な場合：巻数を減らしてまばらにすると、磁場は弱くなり、やがて巻いた導線1本1本の磁場が見えてきます。
\kotae 約 $2.5\times10^{-3}$\,T
\end{phreidai}
\begin{jibun}
長さ\qty{20}{cm}に400回巻いたソレノイドに、（1）\qty{0.50}{A}を流したときの中の磁場を求めましょう。（2）中の磁場を\qty{0.010}{T}にするには、何 A 流せばよいですか。
\bkjans{（1）$n=2000$ /m なので、$4\pi\times10^{-7}\times2000\times0.50$ で約 $1.3\times10^{-3}$\,T。（2）$0.010\div(4\pi\times10^{-7}\times2000)$ で約\qty{4.0}{A}}
\end{jibun}

\label{ph:電-6:ex4}
\begin{phreidai}[環状コイルの中の磁場]
500回巻いた環状コイルに\qty{1.0}{A}の電流を流します。ドーナツの中心線（半径\qty{0.10}{m}）の上の磁場を求めましょう。また、ドーナツの内のはし（半径\qty{0.09}{m}）と外のはし（\qty{0.11}{m}）では、どれだけ違いますか。
\phbamen
ドーナツの形のコイルの中の磁場が、場所でどう違うかを調べる場面です。
\phmushi
巻き方は一様とします。導線の太さは無視します。
\phzu
図\ref{ph:電-6:fig-toroid}。$N=500$、$I=\qty{1.0}{A}$、$r=0.09$・$0.10$・$\qty{0.11}{m}$。
\phriyu
中心が同じ円の道を選ぶと、アンペールの法則から $B\times2\pi r=\mu_0NI$（⑥）。
\phtateru
\[B=\frac{\mu_0NI}{2\pi r}=\frac{2\times10^{-7}\times NI}{r}\]
\phtoku
\[r=0.10:\ B=\frac{2\times10^{-7}\times500}{0.10}=1.0\times10^{-3}\ \text{T}\]
\[r=0.09:\ B\approx1.1\times10^{-3}\ \text{T},\qquad r=0.11:\ B\approx0.91\times10^{-3}\ \text{T}\]
\phtashikame
単位：T。\\
大きさ：内のはしは外のはしより約2割強い（$0.11/0.09\approx1.22$）。ドーナツが細い（太さが半径より小さい）ほど、中の磁場はそろいます。\\
極端な場合：ドーナツの半径をとても大きくすると、まっすぐなソレノイドに近づき、$N/(2\pi r)=n$ として⑥の $\mu_0nI$ になります。
\kotae 中心線で $1.0\times10^{-3}$\,T。内のはし約 $1.1\times10^{-3}$\,T、外のはし約 $0.91\times10^{-3}$\,T
\end{phreidai}
\begin{jibun}
1000回巻いた環状コイルに\qty{2.0}{A}を流します。半径\qty{5.0}{cm}の円の上の磁場を求めましょう。
\bkjans{$2\times10^{-7}\times1000\times2.0\div0.050$ で、$8.0\times10^{-3}$\,T}
\end{jibun}

\label{ph:電-6:ex5}
\begin{phreidai}[磁場の中の導線が受ける力]
\qty{0.50}{T}の一様な磁場の中に、長さ\qty{0.20}{m}の導線を置き、\qty{3.0}{A}の電流を流します。（1）導線が磁場に直角なとき、（2）導線と磁場の間の角が30°のとき、導線が受ける力の大きさを求めましょう。
\phbamen
4-14⑧と同じ数で、電流が磁場から受ける力を求める場面です。
\phmushi
電流を運ぶほかの線は、磁場の外にあるとします。
\phzu
⑧の図。$I=\qty{3.0}{A}$、$B=\qty{0.50}{T}$、$l=\qty{0.20}{m}$。
\phriyu
導線の中の電荷が受けるローレンツ力を全部たした $F=IBl\sin\theta$（⑧）。
\phtateru
\[F=IBl\sin\theta\]
\phtoku
（1）$F=3.0\times0.50\times0.20\times\sin90^\circ=\qty{0.30}{N}$\quad（2）$0.30\times\sin30^\circ=\qty{0.15}{N}$
\phtashikame
単位：A×T×m＝A×N/(A·m)×m＝N。\\
大きさ：\qty{0.30}{N}は、\qty{30}{g}のおもりの重さくらいです。4-14⑧の答えと同じです。\\
極端な場合：導線を磁場と平行にすると（$\theta=0$）、力は0です。
\kotae （1）\qty{0.30}{N}　（2）\qty{0.15}{N}
\end{phreidai}
\begin{jibun}
\qty{0.40}{T}の磁場の中の長さ\qty{0.10}{m}の導線に\qty{2.0}{A}を流します。導線と磁場の角が（1）90°、（2）60°のときの力を求めましょう。
\bkjans{（1）\qty{0.080}{N}。（2）$0.080\times\sin60^\circ$ で約\qty{0.069}{N}}
\end{jibun}

\label{ph:電-6:ex6}
\begin{phreidai}[平行電流の力とアンペアの定義]
真空の中で\qty{1.0}{m}離して平行に置いた長い2本の導線に、同じ向きに\qty{1.0}{A}ずつ電流を流します。導線 1 m あたりにはたらく力の大きさと向きを求めましょう。
\phbamen
⑨の、2019年までのアンペアの定義の場面です。
\phmushi
導線は無限に長く、とても細いとします。
\phzu
図\ref{ph:電-6:fig-parallel}の左。$I_1=I_2=\qty{1.0}{A}$、$r=\qty{1.0}{m}$、$l=\qty{1.0}{m}$。
\phriyu
導線1の磁場 $\mu_0I_1/(2\pi r)$ の中で、導線2が力 $I_2Bl$ を受ける（⑨）。
\phtateru
\[F=\frac{\mu_0I_1I_2l}{2\pi r}\]
\phtoku
\[F=\frac{4\pi\times10^{-7}\times1.0\times1.0\times1.0}{2\pi\times1.0}=2.0\times10^{-7}\ \text{N}\]
同じ向きの電流なので、引き合います。
\phtashikame
単位：N/A$^2$×A×A×m÷m＝N。\\
大きさ：とても小さな力です。\qty{1}{A}は、ふつうの電気製品を流れる電流くらいですが、離れた導線どうしの磁場の力は弱いのです。\\
極端な場合：アンペアの昔の定義の数（$2\times10^{-7}$\,N）がそのまま出てきました。
\kotae 1 m あたり $2.0\times10^{-7}$\,N、引き合う
\end{phreidai}
\begin{jibun}
\qty{2.0}{cm}離した平行な導線に、同じ向きに\qty{5.0}{A}と\qty{3.0}{A}の電流を流します。1 m あたりの力の大きさと向きを求めましょう。
\bkjans{$2\times10^{-7}\times5.0\times3.0\div0.020$ で、$1.5\times10^{-4}$\,N。引き合う}
\end{jibun}

\label{ph:電-6:ex7}
\begin{phreidai}[電気のコードの2本の線の力]
ドライヤーのコードの中で、行きと帰りの2本の線が\qty{3.0}{mm}離れて並び、\qty{10}{A}の電流が流れています。1 m あたりの力を求めましょう。ショート（短絡）して、一瞬\qty{1000}{A}の電流が流れたとすると、力は何倍になりますか。
\phbamen
身近なコードで、平行電流の力の大きさを見積もる場面です。
\phmushi
線は細いまっすぐな導線とします（線の太さ\qty{1}{mm}くらいは、間の\qty{3}{mm}に比べて無視）。交流ですが、行きと帰りの電流はいつも反対向きなので、向きの関係は変わりません（力の大きさは、ある瞬間の\qty{10}{A}で見積もる）。ショートのときの\qty{1000}{A}は、仮の数とします。
\phzu
図\ref{ph:電-6:fig-parallel}の右。$I_1=I_2=\qty{10}{A}$（反対向き）、$r=\qty{0.0030}{m}$。
\phriyu
平行電流の力 $\mu_0I_1I_2l/(2\pi r)$（⑨）。力は電流の2乗に比例する。
\phtateru
\[\frac{F}{l}=\frac{\mu_0I^2}{2\pi r}\]
\phtoku
\[\frac{F}{l}=\frac{2\times10^{-7}\times10^2}{0.0030}\approx6.7\times10^{-3}\ \text{N/m}\]
反対向きの電流なので退け合います。\qty{1000}{A}なら電流は100倍なので、力は $100^2=\num{10000}$ 倍の約\qty{67}{N/m}です。
\phtashikame
単位：N/m。\\
大きさ：ふだんは\qty{7}{mN}（\qty{0.7}{g}のおもりの重さくらい）で、気づかない大きさです。ショートの電流では、1 m あたり\qty{7}{kg}のおもりの重さくらいの力で、線が一瞬はね合います。変電所の太い導体は、この力にたえるように支えます。\\
極端な場合：電流が0なら力は0。
\kotae 約 $6.7\times10^{-3}$\,N/m（退け合う）、\num{10000}倍
\end{phreidai}
\begin{jibun}
2本の線が\qty{4.0}{mm}離れたコードに\qty{15}{A}が流れています。1 m あたりの力を求めましょう。
\bkjans{$2\times10^{-7}\times15^2\div0.0040$ で、約 $1.1\times10^{-2}$\,N/m（退け合う）}
\end{jibun}

\label{ph:電-6:ex8}
\begin{phreidai}[電子が磁場の中で回る半径]
\qty{1.0}{kV}の電圧で加速した電子が、\qty{1.0}{mT}の一様な磁場に直角に入りました。回る円の半径を求めましょう。電子の質量は $9.11\times10^{-31}$\,kg、電荷の大きさは $e=1.60\times10^{-19}$\,C です。
\phbamen
電場で加速した電子が、磁場で曲げられる場面です（ブラウン管や、電子の比電荷をはかる実験）。
\phmushi
重力は無視します（電子の重さは $10^{-29}$\,N の桁で、磁場の力よりずっと小さい）。速さは光の速さよりずっと小さいので、相対論は考えません。はじめの速さは0とします。
\phzu
図\ref{ph:電-6:fig-cyclo}（電子は負なので左回り）。$V=\qty{1000}{V}$、$B=1.0\times10^{-3}$\,T。
\phriyu
加速で電子が得る運動エネルギーは $eV$（電-2 ⑥の電子ボルト）。速さがわかれば、円運動の半径は $r=mv/(eB)$（⑩）。
\phtateru
\[\frac12mv^2=eV,\qquad r=\frac{mv}{eB}\]
\phtoku
\[v=\sqrt{\frac{2eV}{m}}=\sqrt{\frac{2\times1.60\times10^{-19}\times1000}{9.11\times10^{-31}}}\approx1.87\times10^7\ \text{m/s}\]
\[r=\frac{9.11\times10^{-31}\times1.87\times10^7}{1.60\times10^{-19}\times1.0\times10^{-3}}\approx0.11\ \text{m}\]
\phtashikame
単位：C×T＝C×N/(A·m)＝N·s/m なので、kg×(m/s)÷(N·s/m)＝m。\\
大きさ：電-2 例題5 と同じ速さです。光の速さの約6\,\% なので、相対論を無視してよい。半径\qty{11}{cm}の円は、実験の装置の大きさにちょうどよい。\\
極端な場合：$B$ を0にすると $r$ は無限大（まっすぐ進む）。
\kotae 約\qty{0.11}{m}
\end{phreidai}
\begin{jibun}
\qty{100}{V}で加速した電子が\qty{1.0}{mT}の磁場に直角に入ると、半径は何 cm ですか。例題8と比べて、加速の電圧と半径の関係を答えましょう。
\bkjans{$v\approx5.93\times10^6$\,m/s、$r\approx\qty{3.4}{cm}$。電圧が $\tfrac{1}{10}$ で半径は約 $\tfrac{1}{\sqrt{10}}$（半径は電圧の平方根に比例）}
\end{jibun}

\label{ph:電-6:ex9}
\begin{phreidai}[サイクロトロンの周波数]
\qty{1.0}{T}の磁場の中で、陽子（質量 $1.67\times10^{-27}$\,kg、電荷 $1.60\times10^{-19}$\,C）が回る周波数を求めましょう。速さが $1.0\times10^7$\,m/s のときの半径も求めましょう。
\phbamen
サイクロトロンで、加速の電圧の周波数を決める場面です。
\phmushi
相対論は考えません（光の速さの3\,\% ほど）。磁場は一様とします。
\phzu
図\ref{ph:電-6:fig-cyclo}。$B=\qty{1.0}{T}$、$m=1.67\times10^{-27}$\,kg、$q=1.60\times10^{-19}$\,C。
\phriyu
円運動の周波数 $f=qB/(2\pi m)$（⑩）は速さによらない。半径は $r=mv/(qB)$。
\phtateru
\[f=\frac{qB}{2\pi m},\qquad r=\frac{mv}{qB}\]
\phtoku
\[f=\frac{1.60\times10^{-19}\times1.0}{2\pi\times1.67\times10^{-27}}\approx1.5\times10^7\ \text{Hz}=\qty{15}{MHz}\]
\[r=\frac{1.67\times10^{-27}\times1.0\times10^7}{1.60\times10^{-19}\times1.0}\approx0.10\ \text{m}\]
\phtashikame
単位：C×T÷kg＝C×N/(A·m)÷kg＝1/s。\\
大きさ：\qty{15}{MHz}は、短波のラジオくらいの周波数で、電気の回路でつくりやすい周波数です。\\
極端な場合：粒子が重いほど、ゆっくり回ります。
\kotae 約\qty{15}{MHz}、半径約\qty{0.10}{m}
\end{phreidai}
\begin{jibun}
電子（質量 $9.11\times10^{-31}$\,kg）が\qty{1.0}{mT}の磁場の中で回る周波数を求めましょう。
\bkjans{$f=qB/(2\pi m)$ に、$q=1.60\times10^{-19}$、$B=1.0\times10^{-3}$、$m=9.11\times10^{-31}$ を入れて、約\qty{28}{MHz}}
\end{jibun}

\label{ph:電-6:ex10}
\begin{phreidai}[ホール素子で磁場をはかる]
厚さ\qty{0.20}{mm}の半導体のホール素子に、\qty{1.0}{mA}の電流を流し、\qty{0.10}{T}の磁場を垂直にかけます。電流を運ぶ電荷（電気素量 $1.60\times10^{-19}$\,C）の数密度を 1 m$^3$ あたり $1.0\times10^{22}$ 個とします。ホール電圧を求めましょう。
\phbamen
ホール素子の出す電圧から、磁場をはかるしくみを数で確かめる場面です。
\phmushi
電荷はみな同じ速さで流れるとします。素子の温度による変化は無視します。
\phzu
⑩のホール効果の図。$I=1.0\times10^{-3}$\,A、$B=\qty{0.10}{T}$、$n_{\mathrm c}=1.0\times10^{22}$ /m$^3$、$t_{\mathrm H}=2.0\times10^{-4}$\,m。
\phriyu
横の電場の力とローレンツ力のつり合いから、$V_{\mathrm H}=IB/(n_{\mathrm c}qt_{\mathrm H})$（⑩）。
\phtateru
\[V_{\mathrm H}=\frac{IB}{n_{\mathrm c}qt_{\mathrm H}}\]
\phtoku
\[V_{\mathrm H}=\frac{1.0\times10^{-3}\times0.10}{1.0\times10^{22}\times1.60\times10^{-19}\times2.0\times10^{-4}}=\frac{1.0\times10^{-4}}{0.32}\approx3.1\times10^{-4}\ \text{V}\]
\phtashikame
単位：A×T÷(1/m$^3$×C×m)＝A×T×m$^2$÷C＝T×m$^2$÷s。T×m$^2$＝N·m/A なので、N·m/(A·s)＝J/C＝V。\\
大きさ：\qty{0.31}{mV}は、増幅器で十分はかれる大きさです。銅のように数密度が大きい（$10^{29}$ /m$^3$ の桁）と、電圧は数百万分の1になり、はかりにくくなります（章末問題）。\\
極端な場合：$B=0$ なら電圧は0。$B$ の向きが反対なら、電圧の符号が反対。
\kotae 約 $3.1\times10^{-4}$\,V（\qty{0.31}{mV}）
\end{phreidai}
\begin{jibun}
例題10の素子で、ホール電圧が\qty{0.50}{mV}でした。磁場の強さを求めましょう。
\bkjans{$B=V_{\mathrm H}n_{\mathrm c}qt_{\mathrm H}/I$。$n_{\mathrm c}qt_{\mathrm H}=0.32$ なので、$0.50\times10^{-3}\times0.32\div(1.0\times10^{-3})$ で、\qty{0.16}{T}}
\end{jibun}

\label{ph:電-6:ex11}
\begin{phreidai}[電流計のコイルにはたらくトルク]
1辺\qty{2.0}{cm}の正方形のコイルを100回巻き、\qty{0.20}{T}の磁場の中に、コイルの面が磁場と平行になるように置きます。\qty{10}{mA}の電流を流したときのトルクを求めましょう。
\phbamen
アナログの電流計の針を回すトルクの大きさを求める場面です。
\phmushi
磁場は一様とします（本物の電流計では、磁石の形をくふうして、コイルの辺がいつも磁場と直角になるようにしてある）。
\phzu
図\ref{ph:電-6:fig-torque}。$N=100$、$S=(0.020)^2=4.0\times10^{-4}$\,m$^2$、$I=0.010$\,A、$B=\qty{0.20}{T}$、$\theta=90^\circ$。
\phriyu
コイルにはたらくトルクは $T=NISB\sin\theta$（⑪）。面が磁場と平行なので、面に垂直な向きは磁場と直角。
\phtateru
\[T=NISB\sin\theta\]
\phtoku
\[T=100\times0.010\times4.0\times10^{-4}\times0.20\times1=8.0\times10^{-5}\ \text{N·m}\]
\phtashikame
単位：A×m$^2$×T＝A×m$^2$×N/(A·m)＝N·m。\\
大きさ：うで\qty{1}{cm}の所を、約\qty{0.8}{g}のおもりの重さで押すくらいの、小さなトルクです。細いばねで受けとめると、針が目に見える角度だけ回ります。\\
極端な場合：コイルの面が磁場に直角（$\theta=0$）なら0です。
\kotae $8.0\times10^{-5}$\,N·m
\end{phreidai}
\begin{jibun}
\qty{3.0}{cm}×\qty{4.0}{cm}の長方形のコイルを50回巻き、\qty{0.10}{T}の磁場の中で\qty{20}{mA}を流します。（1）面に垂直な向きと磁場の角が90°のとき、（2）30°のときのトルクを求めましょう。
\bkjans{（1）$50\times0.020\times1.2\times10^{-3}\times0.10$ で、$1.2\times10^{-4}$\,N·m。（2）その $\sin30^\circ=0.5$ 倍で $6.0\times10^{-5}$\,N·m}
\end{jibun}

\label{ph:電-6:ex12}
\begin{phreidai}[送電線の下の磁場]
（1）\qty{1000}{A}の電流が流れる長い電線の、真下\qty{20}{m}の所の磁場を求め、地球の磁場（約 $4.6\times10^{-5}$\,T）と、法令の基準（\qty{200}{\micro T}）と比べましょう。（2）\qty{1000}{A}の行きの線と帰りの線が\qty{1.0}{m}離れて水平に並んでいるとき、2本のまん中の真下\qty{20}{m}の所の磁場はどうなりますか。
\phbamen
送電線の下で、電流のつくる磁場がどれくらいかを見積もる場面です。
\phmushi
電線はまっすぐで十分長いとします。電線のたるみ、地面の影響は無視します。本物の送電線は三相交流（電-11）で、3本の電流の向きの関係がもっと複雑ですが、ここでは1本と、行きと帰りの2本の場合を考えます。電流はある瞬間の値で考えます。
\phzu
（1）1本の電線と、真下の点。（2）電線Aと電線B（電流は反対向き）、まん中の真下の点P。それぞれの電線から点Pまでの距離は $\sqrt{20^2+0.5^2}\approx\qty{20}{m}$。
\phriyu
1本の磁場は $\mu_0I/(2\pi r)$（⑤）。2本のときは、それぞれの磁場をベクトルとしてたす（重ね合わせ）。
\phtateru
\[B=\frac{\mu_0I}{2\pi r}\]
（2）2本の磁場は、それぞれの電線と点Pを結ぶ線に垂直です。電流が反対向きなので、2つの磁場はほとんど反対向きで、2つの線の間の小さな角 $\alpha$（$\tan\tfrac{\alpha}{2}=0.5/20$）の分だけ打ち消し残ります。残りは $B_{\text{残}}\approx B\times\alpha$ です。
\phtoku
（1）\[B=\frac{2\times10^{-7}\times1000}{20}=1.0\times10^{-5}\ \text{T}=\qty{10}{\micro T}\]
地球の磁場の約5分の1、法令の基準の20分の1です。\\
（2）$\alpha\approx2\times0.5/20=0.050$ rad なので、
\[B_{\text{残}}\approx1.0\times10^{-5}\times0.050=5.0\times10^{-7}\ \text{T}=\qty{0.50}{\micro T}\]
\phtashikame
単位：T。\\
大きさ：（2）は（1）の20分の1です。2本の間かく（\qty{1}{m}）が距離（\qty{20}{m}）に比べて小さいほど、よく打ち消し合います。Python で2本の磁場をベクトルとしてたしても、$5.0\times10^{-7}$\,T になりました。\\
極端な場合：2本をぴったり重ねると（間かく0）、磁場は0です（家のコード：②の身近な例）。
\kotae （1）\qty{10}{\micro T}（地球の約5分の1、基準の20分の1）　（2）約\qty{0.50}{\micro T}
\end{phreidai}
\begin{jibun}
\qty{200}{A}が流れる配電線の真下\qty{10}{m}の所の磁場を、（1）1本のとき、（2）反対向きの電流の線が\qty{0.60}{m}離れて並ぶときのまん中の真下で、求めましょう。
\bkjans{（1）$2\times10^{-7}\times200\div10$ で、\qty{4.0}{\micro T}。（2）$\alpha\approx0.60\div10=0.060$ rad で、約\qty{0.24}{\micro T}}
\end{jibun}

\label{ph:電-6:ex13}
\begin{phreidai}[速度選別器（入試の形）]
たがいに直角な電場 $E=1.0\times10^4$\,V/m と磁場 $B=\qty{0.010}{T}$ がかかった所へ、両方に直角に、いろいろな速さの正のイオンを入れます。まっすぐ進むイオンの速さを求めましょう。それより速いイオンは、どちらへ曲がりますか。
\phbamen
電場の力と磁場の力のつり合いで、決まった速さの粒子だけを選ぶ場面です。
\phmushi
重力は無視します。電場・磁場は一様で、はしの乱れは考えません。
\phzu
イオンは右へ進み、電場は下向き（電場の力は下向き $qE$）、磁場はおく向き（ローレンツ力は上向き $qvB$）とします。
\phriyu
まっすぐ進むのは、2つの力がつり合う $qE=qvB$ のとき（⑩）。
\phtateru
\[qE=qvB\qquad\Longrightarrow\qquad v=\frac{E}{B}\]
\phtoku
\[v=\frac{1.0\times10^4}{0.010}=1.0\times10^6\ \text{m/s}\]
これより速いイオンは、ローレンツ力 $qvB$ が電場の力 $qE$ より大きいので、上（ローレンツ力の向き）へ曲がります。
\phtashikame
単位：(V/m)÷T＝(N/C)÷(N/(A·m))＝A·m/C＝m/s。\\
大きさ：電荷 $q$ も質量も式に入らないので、どんなイオンでも、この速さのものだけが通ります。\\
極端な場合：$B=0$ なら、どの速さのイオンも電場で下へ曲がります。
\kotae $1.0\times10^6$\,m/s。速いイオンは磁場の力の向き（上）へ曲がる
\end{phreidai}
\begin{jibun}
電場 $2.0\times10^4$\,V/m の速度選別器で、$4.0\times10^6$\,m/s のイオンだけを通すには、磁場を何 T にすればよいですか。
\bkjans{$B=E/v=2.0\times10^4\div(4.0\times10^6)$ で、$5.0\times10^{-3}$\,T}
\end{jibun}

%% ======================== まちがえやすいところ ========================
\Needspace{11\baselineskip}\bksection{tsumazuki}{まちがえやすいところをチェック}
\begin{tsumazuki}{磁場の力の向きは、磁場の向き？}
\bkng{電気の力は電場の向きにはたらくので、磁場の力も磁場の向きにはたらく。}
\bkok{磁場の力（ローレンツ力）は、速度にも磁場にも\textbf{垂直}な向き。}
外積 $\vec v\times\vec B$ の向きを、右ねじかフレミングの左手で調べます（③）。止まった電荷には、磁場の力ははたらきません。
\end{tsumazuki}

\begin{tsumazuki}{電子の流れで右ねじを使う}
\bkng{導線の中を電子が北へ流れているので、北へ進む右ねじで磁場の向きを決める。}
\bkok{電子が北へ流れるとき、電流は南へ流れている。右ねじは\textbf{電流の向き}に進める。}
電流の向きは、正の電荷が流れる向きと決められています（電-4 ①）。ホール効果やローレンツ力の向きも、電子では反対になります（⑩）。
\end{tsumazuki}

\begin{tsumazuki}{まっすぐな電流の磁場も、距離の2乗に反比例？}
\bkng{点電荷の電場と同じで、距離が2倍なら $\tfrac14$。}
\bkok{長いまっすぐな電流の磁場は $1/r$ に比例。距離が2倍なら半分。}
ビオ・サバールの法則の1つ1つの部分は $1/r^2$ ですが、長い導線全体をたすと $1/r$ になります（⑤）。円い電流の中心は、ふつう半径で書くので注意します（$\mu_0I/(2a)$ と $\mu_0I/(2\pi r)$ を取り違えない）。
\end{tsumazuki}

\begin{tsumazuki}{同じ向きの電流は退け合う？}
\bkng{同じ符号の電荷が退け合うので、同じ向きの電流も退け合う。}
\bkok{\textbf{同じ向きの電流は引き合い}、反対向きは退け合う。}
電荷の力と、電流の力は、しくみが違います（⑨）。右ねじと外積で、1つずつ確かめましょう。
\end{tsumazuki}

\begin{tsumazuki}{ローレンツ力で、粒子は速くなる？}
\bkng{磁場の力で押されるので、粒子は加速されて速くなる。}
\bkok{ローレンツ力は速度にいつも垂直なので仕事をせず、速さは変わらない。向きだけが変わる。}
サイクロトロンで粒子を速くしているのは、すき間にかけた\textbf{電場}です（⑩）。
\end{tsumazuki}

%% ======================== 工学ではこう使う ========================
\Needspace{18\baselineskip}\bksection{phgear}{工学ではこう使う}
\begin{phkougaku}[電流と磁場を、はかる・使う]
\textbf{電流センサー}：導線を切らずに電流をはかるクランプメーターや、電気自動車の電池の電流センサーでは、鉄の輪で導線を囲み、輪のすき間にホール素子を置きます。導線の電流がつくる磁場（⑤）を鉄の輪が集め、ホール素子が電圧に変えます。アンペールの法則から、輪の中のどこを導線が通っても、輪を1周する磁場の線積分は電流で決まるので、導線の位置にあまりよりません。

\textbf{MRI}：病院の MRI は、超伝導の線を巻いた大きなソレノイドで、1.5〜\qty{3}{T}の一様な磁場をつくります。⑥の $\mu_0nI$ で\qty{3}{T}をつくるには、$nI\approx2.4\times10^6$\,A/m、たとえば 1 m あたり2000回巻いて約\qty{1200}{A}を流すことになります。ふつうの銅線では熱が出すぎるので、電気抵抗が0になる超伝導の線を使います（超伝導のしくみは物理シリーズの『原子』（原-8））。

\textbf{リニアモーターカーと電磁ブレーキ}：電磁石と、⑧の電流が受ける力を、まっすぐな線路の上に広げたものがリニアモーターです。

\begin{phketa}[電流の磁場・力の目安]
\phketaline{家のコードの電流（10 A）から 5 cm}{$4\times10^{-5}$\,T（地球の磁場と同じくらい）}
\phketaline{ソレノイド（1000 回/m・数 A）}{数 mT}
\phketaline{MRI（超伝導のソレノイド）}{1.5〜\qty{3}{T}}
\phboxhead{答えが桁で合っているか}
$\mu_0/(2\pi)=2\times10^{-7}$ を覚えておくと、まっすぐな電流の磁場は、電流に $2\times10^{-7}$ をかけて距離でわるだけで、すぐ見積もれます。鉄を使わない電流の磁場が\qty{1}{T}をこえるのは、何千 A を流す特別な装置だけです。
\end{phketa}
\end{phkougaku}
\phstation{電気機器,計測工学,医用画像,加速器・プラズマ}{磁束密度、磁界の強さ、起磁力・アンペアターン、ホール素子、電流センサー、ローレンツ力}

\begin{column}{エルステッドの発見と、その年の秋}
エルステッドは、講義の準備か実験の最中に、電流を流した導線の近くで方位磁針が動くのを見たと伝えられています。1820年7月、その結果をラテン語の4ページの報告にまとめ、ヨーロッパの研究者たちに送りました。報告がパリに届くと、アンペールはすぐに実験を始め、9月には、平行な電流が引き合ったり退け合ったりすることを示しました。10月には、ビオとサバールが、導線からの距離と磁場の強さの関係を発表しました。エルステッドの発見から数か月のうちに、この章の②・④・⑨の考えの土台ができたのです。翌1821年には、イギリスのファラデーが、電流と磁石で針金をぐるぐる回し続ける装置をつくりました。モーターの始まりです。
\end{column}

\begin{column}{スマホは、どうやって向きを知る？}
スマホの地図アプリは、持っている人がどちらを向いているかを、画面の矢印で示します。スマホの中には、小さな磁気センサー（電子コンパス）が入っていて、地球の磁場を、たて・横・高さの3つの向きに分けてはかっています（⑩のホール素子や、磁場で電気抵抗が変わる素子）。3つの成分から、⑫の北の向きと伏角がわかります。地球の磁場は $5\times10^{-5}$\,T くらいと弱いので、スマホの中のスピーカーの磁石や、近くの鉄・電流の磁場にじゃまされます。アプリが「スマホを8の字に回してください」と言うのは、いろいろな向きにしたときのはかった値から、スマホ自身がつくる磁場の分を求めて、差し引くためです。偏角を直して真北を出すのは、アプリが位置（GPS）から計算しています。
\end{column}


%% ======================== 章末問題 ========================
\bkendproblems
\bkproblemgroup{基本}{3}
\bkproblem 長いまっすぐな導線に\qty{20}{A}の電流が流れています。導線から（1）\qty{4.0}{cm}、（2）\qty{8.0}{cm}の所の磁場を求めましょう。
\bkpans{（1）$2\times10^{-7}\times20\div0.040$ で、$1.0\times10^{-4}$\,T。（2）距離が2倍なので半分の $5.0\times10^{-5}$\,T}
\bkproblem 長さ\qty{0.30}{m}に600回巻いたソレノイドに\qty{1.5}{A}を流します。中の磁場を求めましょう。
\bkpans{$n=2000$ /m。$4\pi\times10^{-7}\times2000\times1.5$ で、約 $3.8\times10^{-3}$\,T}
\bkproblem 半径\qty{4.0}{cm}、20回巻きの円いコイルに\qty{0.80}{A}を流します。中心の磁場を求めましょう。
\bkpans{$4\pi\times10^{-7}\times20\times0.80\div(2\times0.040)$ で、約 $2.5\times10^{-4}$\,T}
\bkproblem \qty{0.30}{T}の一様な磁場の中に、長さ\qty{0.50}{m}の導線を置き、\qty{4.0}{A}を流します。導線が（1）磁場と直角のとき、（2）磁場と平行のときの力の大きさを求めましょう。
\bkpans{（1）$4.0\times0.30\times0.50$ で、\qty{0.60}{N}。（2）0}
\bkproblem \qty{5.0}{cm}離した平行な長い導線に、反対向きに\qty{10}{A}ずつ流します。1 m あたりの力の大きさと向きを求めましょう。
\bkpans{$2\times10^{-7}\times10\times10\div0.050$ で、$4.0\times10^{-4}$\,N。退け合う}
\bkproblem 陽子（質量 $1.67\times10^{-27}$\,kg、電荷 $1.60\times10^{-19}$\,C）が、速さ $2.0\times10^6$\,m/s で\qty{0.50}{T}の磁場に直角に入りました。回る円の半径を求めましょう。
\bkpans{$r=mv/(qB)$ に、$m=1.67\times10^{-27}$、$v=2.0\times10^6$、$q=1.60\times10^{-19}$、$B=0.50$ を入れて、約\qty{4.2}{cm}}
\bkproblem 面積\qty{5.0}{cm^2}のコイルを200回巻き、\qty{0.30}{T}の磁場の中に、面が磁場と平行になるように置いて\qty{50}{mA}を流します。トルクを求めましょう。
\bkpans{$T=NISB=200\times0.050\times5.0\times10^{-4}\times0.30$ で、$1.5\times10^{-3}$\,N·m}
\bkproblem アルファ粒子（ヘリウムの原子核。質量 $6.64\times10^{-27}$\,kg、電荷は陽子の2倍の $3.20\times10^{-19}$\,C）が\qty{1.0}{T}の磁場の中で回る周波数を求めましょう。
\bkpans{$f=qB/(2\pi m)$ に、$q=3.20\times10^{-19}$、$B=1.0$、$m=6.64\times10^{-27}$ を入れて、約\qty{7.7}{MHz}}
\bkproblemgroup{チャレンジ}{4}
\bkproblem 質量分析器：速度選別器（$E=2.0\times10^4$\,V/m、$B=\qty{0.050}{T}$）を通りぬけた正のイオンを、\qty{0.050}{T}の一様な磁場に直角に入れ、半円を回らせて、入った所から直径の所の板に当てます。陽子（質量 $1.67\times10^{-27}$\,kg）と、質量が2倍の重水素の原子核（電荷は同じ $1.60\times10^{-19}$\,C）では、板に当たる点は何 cm 離れますか（例題8・例題13）。
\bkpans{選ばれる速さ $v=E/B=4.0\times10^5$\,m/s。陽子の半径 $r=mv/(qB)\approx\qty{8.4}{cm}$、重水素は2倍の約\qty{16.7}{cm}。当たる点は入口から直径の所なので、差は $2\times(16.7-8.4)$ で約\qty{17}{cm}}
\bkproblem 反対向きに\qty{10}{A}ずつ流れる平行な長い導線が\qty{10}{cm}離れています。2本のまん中の点の磁場を求めましょう。同じ向きに流したときはどうですか（例題12）。
\bkpans{反対向き：2本の磁場はまん中で同じ向きになり、$2\times(2\times10^{-7}\times10\div0.050)$ で、$8.0\times10^{-5}$\,T。同じ向き：反対向きで打ち消し合い、0}
\bkproblem 内半径\qty{8.0}{cm}、外半径\qty{12}{cm}のドーナツに400回巻いた環状コイルに\qty{2.0}{A}を流します。内のはしと外のはしの磁場と、その比を求めましょう（例題4）。
\bkpans{内：$2\times10^{-7}\times800\div0.080$ で $2.0\times10^{-3}$\,T。外：$\div0.12$ で約 $1.3\times10^{-3}$\,T。比は $12\div8=1.5$}
\bkproblem 銅の板（厚さ\qty{0.10}{mm}、自由電子の数密度 $8.5\times10^{28}$ /m$^3$）に\qty{10}{A}を流し、\qty{1.0}{T}の磁場を垂直にかけます。ホール電圧を求め、例題10の半導体の素子と比べて、なぜホール素子に半導体を使うのかを説明しましょう。
\bkpans{$V_{\mathrm H}=IB/(n_{\mathrm c}qt_{\mathrm H})$ に、$I=10$、$B=1.0$、$n_{\mathrm c}=8.5\times10^{28}$、$q=1.60\times10^{-19}$、$t_{\mathrm H}=1.0\times10^{-4}$ を入れて、約 $7.4\times10^{-6}$\,V。大きな電流と強い磁場でも\qty{7}{\micro V}ほどしか出ない。数密度が小さい半導体のほうが、同じ電流・磁場で大きな電圧が出るから}

\begin{furikaeri}
\begin{checklist}
\item 電流の周りに輪の形の磁場ができることを方位磁針で確かめ、その向きを右ねじで言える。\bkfurref{ph:電-6:sec2}
\item まっすぐな電流・円い電流・ソレノイドがつくる磁場と、磁場の中の電流・荷電粒子・コイルが受ける力を求められる。\bkfurref{ph:電-6:sec4}
\item ビオ・サバールの法則とアンペールの法則（積分形）から磁場を導き、$\mathrm{rot}\,\vec B=\mu_0\vec j$ と $\mathrm{div}\,\vec B=0$ を読める。\bkfurref{ph:電-6:sec6}
\item ホール効果・速度選別器・サイクロトロンのしくみを、ローレンツ力で説明できる。\bkfurref{ph:電-6:sec10}
\end{checklist}
\end{furikaeri}
\bkfukidashi{sora}{導線の電流が、地球と同じくらいの磁場をつくっていたから、方位磁針が大きくふれたんだね。モーターもスピーカーも、同じ $q\vec v\times\vec B$ で動いているんだ！}
